\documentclass[10pt,a4paper]{article}

% ----------------- Paquetes -------------------%
\usepackage[T1]{fontenc}
\usepackage[utf8]{inputenc}
\usepackage[a4paper,margin=2.5cm]{geometry}
\usepackage{mathptmx}             
\usepackage{changepage} 
\usepackage{setspace}
\usepackage[spanish]{babel}
\usepackage[labelfont=bf,labelsep=space]{caption}
\usepackage{amssymb}
\usepackage{amsmath}
\usepackage{ebgaramond}
\usepackage{placeins}
\usepackage{etoolbox}
\usepackage{setspace}
\usepackage{parskip} 
\usepackage{tcolorbox}
\usepackage{needspace}
\usepackage{xparse}
\usepackage{ocgx2}
\usepackage{hyperref}
\usepackage{hypcap}
\usepackage{soul}\sethlcolor{yellow}% resaltado amarillo: \hl{texto} (Panel TCEE, ⌃H)



%%%%%%%%%%%%%%%%%%%%%%%%%%%%%%%%%%%%%%%%%%%%%%%%%%%%%%%%%%%%%%%%
% --- Fija primero tus valores globales "buenos" ---
\setstretch{1.2}                 % Interlineado global
\setlength{\parskip}{0.7\baselineskip}
\setlength{\parindent}{0pt}
\newlength{\GlobalParskip}
\setlength{\GlobalParskip}{\parskip}

\newcounter{eqnum}[subsection]
\renewcommand{\theeqnum}{\arabic{eqnum}}
\newcounter{alphaitem}
\newcounter{numitem}

%% ------------------- Recuadros----------------------%%
\newcommand{\puntosclave}[1]{%
  \begin{tcolorbox}[
    colframe=black,
    title={Puntos clave},
    before upper={\setlength{\parindent}{0.5cm}\setlength{\parskip}{0.5\baselineskip}}
  ]
  #1
  \end{tcolorbox}
}

\newcommand{\modificaciones}[1]{
  \begin{tcolorbox}[
    colback=white,
    colframe=red,
    title={Modificaciones a realizar},
    before upper={\setlength{\parindent}{0.5cm}\setlength{\parskip}{0.5\baselineskip}}
  ]
  #1
  \end{tcolorbox}
}

%% ------------------- Preguntas TEST----------------------%%
% Opciones: radio (single-choice)
\NewDocumentCommand{\OptRadio}{m m m}{%
  \noindent\CheckBox[radio,name=#1,value=#2,width=1.2ex,height=1.2ex]{}%
  \;\textbf{#2.}~#3\par
}

% Opciones: checkbox (multi-choice)
\NewDocumentCommand{\OptCheck}{m m}{%
  \noindent\CheckBox[name=#1,width=1.2ex,height=1.2ex,bordercolor={0 0 0}]{}%
  \;#2\par
}

% Pregunta single-choice (radio)
% Args: 1=id, 2=enunciado, 3=opciones, 4=solución+justificación, 5=imagen (o {}), 6=origen
\NewDocumentCommand{\PreguntaTestSingle}{m m m m m m}{%
  \par\medskip
  \noindent\textbf{#1.}~#2\par\smallskip

    % Imagen (si no es {})
    \IfBlankTF{#5}{}{%
      \begin{center}
        \includegraphics[width=0.75\linewidth]{#5}
      \end{center}
    }

    \begin{Form}
      #3
    \end{Form}

    \par\smallskip
    \switchocg{sol-#1}{\fbox{Mostrar / ocultar solución}}%
    \hfill{\footnotesize #6}\par\smallskip

    \begin{ocg}{Solución}{sol-#1}{0}
      \noindent #4
    \end{ocg}
  \par\medskip
}

% Pregunta multi-choice (checkboxes)
\NewDocumentCommand{\PreguntaTestMulti}{m m m m m m}{%
  \par\medskip
  \noindent\textbf{#1.}~#2\par\smallskip

    % Imagen (si no es {})
    \IfBlankTF{#5}{}{%
      \begin{center}
        \includegraphics[width=0.75\linewidth]{#5}
      \end{center}
    }
    \begin{Form}
      #3
    \end{Form}

    \par\smallskip
    \switchocg{sol-#1}{\fbox{Mostrar / ocultar solución}}%
    \hfill{\footnotesize #6}\par\smallskip

    \begin{ocg}{Solución}{sol-#1}{0}
      \noindent #4
    \end{ocg}
  \par\medskip
}

%% ------------------- CITA ----------------------%%
% \begin{cita}[Autor][Año][Obra] texto \end{cita}: cita sangrada y, debajo a la derecha, «AUTOR (Año), Obra». Los tres datos son opcionales.
% En el Panel TCEE: escribir \cita e Intro (el cursor va al texto; Tab pasa a Autor, Año y Obra).
\NewDocumentEnvironment{cita}{ O{} O{} O{} +b }{%
  \begin{adjustwidth}{2cm}{0pt}
  \raggedright
  #4\par
  \raggedleft
  \if\relax\detokenize{#1#2#3}\relax
    % sin firma
  \else
    #1%
    \if\relax\detokenize{#2}\relax\else\if\relax\detokenize{#1}\relax\else\space\fi(#2)\fi
    \if\relax\detokenize{#3}\relax\else\if\relax\detokenize{#1#2}\relax\else, \fi\textit{#3}\fi
  \fi
  \end{adjustwidth}
}{}



%% ------------------- Listas----------------------%%
    % --- Lista alfabética ---
    \newenvironment{la}
      {\begin{list}{\alph{alphaitem})}{%
          \usecounter{alphaitem}%
          \setlength{\leftmargin}{1cm}%
          \setlength{\parskip}{3pt}% 
          \setlength{\itemsep}{0pt}% ← espacio entre ítems
          \setlength{\parsep}{0pt}%  ← párrafos dentro de un ítem
      }}
      {\end{list}}
      
    % --- Lista numérica ---
    \newenvironment{lnum}
      {\begin{list}{\arabic{numitem}.}{%
          \usecounter{numitem}%
          \setlength{\leftmargin}{1cm}%
          \setlength{\parskip}{3pt}% 
          \setlength{\itemsep}{0pt}% ← espacio entre ítems
          \setlength{\parsep}{0pt}%  ← párrafos dentro de un ítem
      }}
      {\end{list}}

%% ------------------- Preguntas Test ----------------------%%
      \usepackage{xparse} % si ya lo tienes, no lo repitas

    \NewDocumentEnvironment{test}{m m o}{%
      \begin{tcolorbox}[
        colback=white,
        colframe=black!15,
        boxrule=0.6pt,
        arc=2mm,
        left=2mm,right=2mm,top=2mm,bottom=2mm
      ]
      \centering
      \includegraphics[width=\linewidth]{#1}
      \end{tcolorbox}
    
      \vspace{2mm}
    
      % Caja SOLO para texto (SÍ partible y mantiene márgenes al partir)
      \begin{tcolorbox}[
        enhanced,
        breakable,
        colback=black!3,
        colframe=black!25,
        boxrule=0.6pt,
        arc=2mm,
        left=2mm,right=2mm,top=1.5mm,bottom=1.5mm
      ]
      \textbf{Respuesta:} \textbf{#2}%
      \IfValueT{#3}{\par\smallskip \textbf{Justificación:} #3}
      \end{tcolorbox}
    }{}
      
% ------------------- Imágenes----------------------%
    \graphicspath{{Img/}}% Rutas por defecto 
    % --- Imagen adaptativa ---
    % Registros de dimensión definidos una sola vez
    \newdimen\imgcap
    \newdimen\imgmin
    \newdimen\imgmax
    \newdimen\imgremain
    \newdimen\imgh
    \newdimen\imgneed
    
    \newcommand{\imagenfit}[3]{%
      \addvspace{10pt}% separación antes
      \par\begingroup
        % Parámetros de composición (asignaciones, no \newdimen)
        \imgcap=1\baselineskip            % alto estimado de título+fuente
        \imgmin=0.15\textheight
        \imgmax=0.15\textheight
        \imgremain=\pagegoal
        \ifdim\pagegoal=\maxdimen \imgremain=0pt\fi
        \advance\imgremain by -\pagetotal
        \advance\imgremain by -\imgcap
        % Decisión de altura destino
        \ifdim\imgremain<\imgmin
          \newpage
          \imgh=0.20\textheight
        \else
          \imgh=\imgremain
          \ifdim\imgh>\imgmax \imgh=\imgmax \fi
        \fi
        % Reservar espacio para evitar cortes del bloque completo
        \imgneed=\imgh \advance\imgneed by \imgcap
        \Needspace{\imgneed}
        % Cuerpo
        \noindent\begin{minipage}{\textwidth}
          \centering
          \captionof{figure}{\textemdash\ #2}\par\vspace{0.3em}% título arriba
          \includegraphics[width=\linewidth,height=\imgh,keepaspectratio]{\detokenize{#1}}\par
          \vspace{0.3em}\small\textit{Fuente: }#3% fuente abajo
        \end{minipage}%
      \endgroup
      \par\addvspace{10pt}%
    }

%------------ Bloque de RECUADRO ESPECIAL -------------%
    \newenvironment{specialblock}[1]{%
      \begin{quote} % crea sangría a izquierda y derecha
      \small % texto más pequeño
      \noindent\textbf{#1}\\[-0.3em] % título en negrita
      \rule{\linewidth}{0.4pt}\\[0.6em]}{
      \end{quote}}
      
%-------------------------- Portada -----------------------%
\newcommand{\oldtitle}[1]{\def\OldTitle{#1}}
\newcommand{\changerationale}[1]{\def\ChangeWhy{#1}}

\newcommand{\changebox}{%
\begin{tcolorbox}[title={Historial de cambios del tema}]
\textbf{Título anterior:} \OldTitle\par
\textbf{Motivo del cambio:} \ChangeWhy
\end{tcolorbox}
}

%-------------------------- Author Font -----------------------%
    \newcommand{\authorfont}[1]{{\fontfamily{EBGaramond-TLF}\selectfont #1}}

%-------------------------- Anexos -----------------------%
    \makeatletter
    \newcounter{anexo}
    \renewcommand{\theanexo}{\Roman{anexo}}
    \newcommand{\anexo}[1]{%
      \clearpage
      \phantomsection
      \refstepcounter{anexo}%
      \noindent\textbf{Anexo \theanexo.- }#1\par\medskip
      \addcontentsline{stc}{section}{Anexo \theanexo.- #1}%
    }
    \makeatother

    %-------------------------- Lista de figuras (personal) -----------------------%
\makeatletter
\newcommand{\MyListOfFigures}{%
  \clearpage
  \phantomsection
  {\centering\bfseries\Large Índice de figuras\par}%
  \medskip
  % Entrada en nuestro índice de contenido (stc)
  \addcontentsline{stc}{section}{Índice de figuras}%
  % Recuperación del listado (sin cabecera propia)
  \begingroup
    \parindent=0pt\parskip=2pt
    \@starttoc{lof}%
  \endgroup
}
\makeatother

%-------------------------- Bibliografía -----------------------%
\makeatletter
% Contador para las referencias
\newcounter{myref}
% Cabecera de la bibliografía, entra en el índice 'stc'
\newcommand{\MyBibliography}{%
  \clearpage
  \phantomsection
  {\centering\bfseries\Large Bibliografía y referencias\par}%
  \medskip
  \addcontentsline{stc}{section}{Bibliografía y referencias}%
  \setcounter{myref}{0}%
}
\newcommand{\MyBibItem}[4]{%
  \refstepcounter{myref}%
  \noindent[\themyref]\ #1.\ \emph{#2}. #3. #4\par
}
\makeatother

%--------- Establecer cuatro niveles de índice --------%
    \setcounter{tocdepth}{4}   
    \setcounter{secnumdepth}{4}% numera hasta \paragraph

%%%%%%%%%%%%%%%%%%%%%%%%%%%%%%%%

\makeatletter

    %--------- Interlineado Global --------%
    \edef\GlobalBaseLineStretch{\baselinestretch}
    
    %--------- Espaciado de seccinones --------%
    \renewcommand\section{\@startsection{section}
      {1}{\z@}
      {2.5ex \@plus 1ex \@minus .2ex} % espacio antes
      {1.5ex \@plus .2ex}             % espacio después
      {\normalfont\Large\bfseries}    % formato del título
    }
    \renewcommand\subsection{\@startsection{subsection}{2}{\z@}
      {3ex \@plus 0.8ex \@minus .2ex}
      {1.5ex \@plus .2ex}
      {\normalfont\large\bfseries}
    }
    \renewcommand\subsubsection{\@startsection{subsubsection}{3}{\z@}
      {3ex \@plus 0.5ex \@minus .2ex}
      {1ex \@plus .2ex}
      {\normalfont\normalsize\bfseries}
    }
    \renewcommand\paragraph{\@startsection{paragraph}{4}{\z@}%
    {-2ex \@plus -0.5ex \@minus -.2ex}% espacio antes
    {0.8ex \@plus .2ex}% espacio después
    {\normalfont\normalsize\itshape}}% ← cursiva en lugar de negrita

    %--------- Ecuaciones --------%   
    % Variables globales para almacenar títulos y notas
    \gdef\leq@current@section{}%
    \gdef\leq@current@subsection{}%
    \gdef\leq@last@printed@section{}%
    \gdef\leq@last@printed@subsection{}%
    
    % Comando para añadir nota (invisible en el cuerpo)
    \newcommand{\eqnote}[1]{%
      \addcontentsline{leq}{leqnote}{#1}%
    }
    
    % Comando para ecuaciones
    \newcommand{\eqblock}[2]{%
      % Verificar si necesitamos imprimir el título de sección
      \ifx\leq@current@section\leq@last@printed@section\else
        \addcontentsline{leq}{leqsection}{\leq@current@section}%
        \global\let\leq@last@printed@section\leq@current@section
        \gdef\leq@last@printed@subsection{}% Reset subsección al cambiar sección
      \fi
      % Verificar si necesitamos imprimir el título de subsección
      \ifx\leq@current@subsection\leq@last@printed@subsection\else
        \ifx\leq@current@subsection\@empty\else
          \addcontentsline{leq}{leqsubsection}{\leq@current@subsection}%
          \global\let\leq@last@printed@subsection\leq@current@subsection
        \fi
      \fi
      % Ecuación
      \refstepcounter{eqnum}%
      \begin{equation*}
        #1\tag{E.\theeqnum}
      \end{equation*}%
      \addcontentsline{leq}{leqt}{[E.\theeqnum] -- #2}%
      \addcontentsline{leq}{leqm}{#1}%
    }
    
    %%% ----- Índice de ecuaciones ----------%%%
    % Tamaño para []
    \newcommand{\leqIndexFont}{\normalsize}
    \newcommand{\leqTitleFont}{\tiny}
    \newcommand{\leqNoteFont}{\tiny}
    % Cabecera de sección 
    \newcommand*\l@leqsection[2]{%
      \addvspace{25pt}% Mayor separación antes de nueva sección
      \noindent\leqIndexFont
      \makebox[\linewidth]{\rule{.38\linewidth}{0.4pt}\ \ \textbf{  \MakeUppercase{#1}}\ \ \rule{.38\linewidth}{0.4pt}}%
      \par\vspace{-10pt}
    }
    % Cabecera de subsección 
    \newcommand*\l@leqsubsection[2]{%
      \par\addvspace{15pt}%
      \noindent\leqIndexFont #1
      \par\vspace{-7pt}
    }
    % Título de ecuación 
    \newcommand*\l@leqt[2]{%
    \par\addvspace{10pt}%
      \par\noindent\leqTitleFont #1\par
    }
    % Miniatura de ecuación
    \newcommand*\l@leqm[2]{%
      \vspace{0.3\baselineskip}%
      \par\scriptsize\noindent\hspace*{1.5em}\(#1\)\par%
    }
    % Nota de ecuación 
    \newcommand*\l@leqnote[2]{%
      \vspace{0.2\baselineskip}%
      \par\noindent\leqNoteFont\hspace*{1.5em}\textbf{Nota.--} #1\par\vspace{0.5\baselineskip}%
    }
    % Generación del índice
    \newcommand{\mylistofequations}{%
      \clearpage
      \phantomsection
      % Título visual de la lista (ajústalo a tu gusto):
      {\centering\bfseries\Large Índice de ecuaciones\par}
      % Entrada en nuestro índice 'stc':
      \addcontentsline{stc}{section}{Índice de ecuaciones}%
      % Llamada a tu lista real (usa el nombre exacto de tu macro):
      \@starttoc{leq}%
    }
    
    %--------- Captura de títulos de sección --------%
    \let\leq@original@section\section
    \let\leq@original@subsection\subsection
    % Flag para saber si estamos en una sección asterisco
    \newif\ifLeqStarSection
    \LeqStarSectionfalse
    
    % Redefinir \section CON CONTROL DE ASTERISCO
    \renewcommand{\section}{%
      \@ifstar{\leq@section@star}{\@dblarg\leq@section@nostar}%
    }
    \def\leq@section@star#1{%
      \global\LeqStarSectiontrue
      \gdef\leq@current@section{#1}%
      \gdef\leq@current@subsection{}%
      \leq@original@section*{#1}%
      \global\LeqStarSectionfalse
    }
    \def\leq@section@nostar[#1]#2{%
      \global\LeqStarSectionfalse
      \gdef\leq@current@section{#2}%
      \gdef\leq@current@subsection{}%
      \leq@original@section[#1]{#2}%
    }
    % Redefinir \subsection
    \renewcommand{\subsection}{%
      \@ifstar{\leq@subsection@star}{\@dblarg\leq@subsection@nostar}%
    }
    \def\leq@subsection@star#1{%
      \global\LeqStarSectiontrue
      \gdef\leq@current@subsection{#1}%
      \leq@original@subsection*{#1}%
      \global\LeqStarSectionfalse
    }
    \def\leq@subsection@nostar[#1]#2{%
      \global\LeqStarSectionfalse
      \gdef\leq@current@subsection{#2}%
      \leq@original@subsection[#1]{#2}%
    }

    % ----- Restauración de valores Globales de estilo ----------%
    % Flags
    \newif\ifInFootnote
    \InFootnotefalse
    \pretocmd{\@footnotetext}{\InFootnotetrue}{}{}
    \apptocmd{\@footnotetext}{\InFootnotefalse}{}{}
    % Profundidad de listas (soporta anidadas)
    \newcount\MyListDepth \MyListDepth=0
    \AtBeginEnvironment{la}{\advance\MyListDepth by 1}
    \AfterEndEnvironment{la}{\advance\MyListDepth by -1}
    \AtBeginEnvironment{lnum}{\advance\MyListDepth by 1}
    \AfterEndEnvironment{lnum}{\advance\MyListDepth by -1}
    \AtBeginEnvironment{itemize}{\advance\MyListDepth by 1}
    \AfterEndEnvironment{itemize}{\advance\MyListDepth by -1}
    \AtBeginEnvironment{enumerate}{\advance\MyListDepth by 1}
    \AfterEndEnvironment{enumerate}{\advance\MyListDepth by -1}
    % Restauración diferida: hazlo al INICIO del próximo párrafo real
    \newif\if@pendingRestore
    \@pendingRestorefalse
    \newcommand{\RequestRestore}{\global\@pendingRestoretrue}
    \everypar{%
      \if@pendingRestore
        \global\@pendingRestorefalse
        \ifInFootnote\else
          \ifnum\MyListDepth=0
            \setlength{\parskip}{\GlobalParskip}%
            \setstretch{\GlobalBaseLineStretch}%
          \fi
        \fi
      \fi
    }
    % Al final de bloques:
    \AfterEndEnvironment{equation}{\RequestRestore}
    \AfterEndEnvironment{equation*}{\RequestRestore}
    \AfterEndEnvironment{la}{\RequestRestore}
    \AfterEndEnvironment{lnum}{\RequestRestore}
    % Al final de macros:
    \apptocmd{\eqblock}{\RequestRestore}{}{}
    \apptocmd{\imagenfit}{\RequestRestore}{}{}
    
\makeatother

% ===================== ToC propio con 4 niveles (único bloque) =====================
\makeatletter

% --- Escritores: añadimos una línea al .stc en cada cabecera numerada ---
% Guardamos las versiones actuales para llamar al comportamiento normal
\let\MyTOC@orig@section\section
\let\MyTOC@orig@subsection\subsection
\let\MyTOC@orig@subsubsection\subsubsection
\let\MyTOC@orig@paragraph\paragraph

% SECTION
\renewcommand{\section}{%
  \@ifstar{\MyTOC@section@star}{\@dblarg\MyTOC@section@nostar}%
}
\def\MyTOC@section@star#1{%
  \MyTOC@orig@section*{#1}% sin entrada al .stc
}
\def\MyTOC@section@nostar[#1]#2{%
  \MyTOC@orig@section[{#1}]{#2}% comportamiento normal (escribe al .toc)
  % Escribimos también nuestra línea al .stc (texto con \numberline)
  \addcontentsline{stc}{section}{\protect\numberline{\thesection}#2}%
}

% SUBSECTION
\renewcommand{\subsection}{%
  \@ifstar{\MyTOC@subsection@star}{\@dblarg\MyTOC@subsection@nostar}%
}
\def\MyTOC@subsection@star#1{%
  \MyTOC@orig@subsection*{#1}%
}
\def\MyTOC@subsection@nostar[#1]#2{%
  \MyTOC@orig@subsection[{#1}]{#2}%
  \addcontentsline{stc}{subsection}{\protect\numberline{\thesubsection}#2}%
}

% SUBSUBSECTION
\renewcommand{\subsubsection}{%
  \@ifstar{\MyTOC@subsubsection@star}{\@dblarg\MyTOC@subsubsection@nostar}%
}
\def\MyTOC@subsubsection@star#1{%
  \MyTOC@orig@subsubsection*{#1}%
}
\def\MyTOC@subsubsection@nostar[#1]#2{%
  \MyTOC@orig@subsubsection[{#1}]{#2}%
  \addcontentsline{stc}{subsubsection}{\protect\numberline{\thesubsubsection}#2}%
}

% PARAGRAPH
\renewcommand{\paragraph}{%
  \@ifstar{\MyTOC@paragraph@star}{\@dblarg\MyTOC@paragraph@nostar}%
}
\def\MyTOC@paragraph@star#1{%
  \MyTOC@orig@paragraph*{#1}%
}
\def\MyTOC@paragraph@nostar[#1]#2{%
  \MyTOC@orig@paragraph[{#1}]{#2}%
  % Si no numeras los \paragraph, cambia \theparagraph por texto plano sin \numberline
  \addcontentsline{stc}{paragraph}{\protect\numberline{\theparagraph}#2}%
}

% --- Impresor del índice propio (lee el .stc) ---
\newcommand{\MyTOC}{%
  \par\bigskip
  {\centering\bfseries\Large Índice de contenido\par}%
  \medskip
  \begingroup
    \parindent=0pt\parskip=2pt
    \@starttoc{stc}%
  \endgroup
}

\makeatother
% ===================== fin ToC propio =====================


%%%%%%%%%%%%%%


% --- Datos del tema ---
\title{Tema 3.A.13 -- Economía de la información.\\
\large  Teoría de la agencia: selección adversa y riesgo moral.}
\author{Marcelo Ortega}
\date{Febrero 2026}
\newcommand{\TemaCode}{3A13}

\oldtitle{Tema 3.A.13. Economía de la información y teoría de la agencia: selección adversa y riesgo moral}
\changerationale{
    Sin cambios.
}

\begin{document}


\maketitle
\changebox

\newpage

% --- Página de resumen y modificaciones ---
\puntosclave{ %% Escribir puntos clave Apuntes CECO
    }
\vspace{1cm}

\modificaciones{
    
    }

\newpage
\MyTOC
\newpage

\mylistofequations
\clearpage

\MyListOfFigures
\clearpage


\pagenumbering{arabic}
\setcounter{page}{1}


\section{Introducción}\label{intro}


\subsection{Contextualización}\label{context}


\subsubsection{Relevancia}\label{importance}


\subsubsection{Historia}\label{history}



\subsubsection{Problemática}\label{problem} 




\subsection{Estructura de la exposición}\label{structure}





\clearpage
\section{Economía de la información}
\puntosclave{
    La economía de la información analiza relaciones contractuales cuando las partes cuentan con distintos conjuntos de información, lo que afecta a sus incentivos para le buen desarrollo de la relación
    
	Se caracteriza por el uso de modelos de equilibrio parcial y la teoría de juegos para analizar contratos entre partes que poseen información privada relevante
    
	La relación de agencia entre Principal y Agente es la herramienta para modelizar estos contratos, relación caracterizada por el conflicto de intereses entre Principal y Agente, así como por la asimetría informativa

	Los diferentes modelos se pueden clasificar atendiendo al momento en que surge la asimetría informativa (antes o después de la firma del contrato, diferenciando selección adversa y señalización frente a riesgo moral) y al objeto de dicha información asimétrica (información oculta o acción oculta)
}

	
\subsection{Economía del Bienestar e Hipótesis}
El Primer Teorema Fundamental de la Economía del Bienestar es uno de los principales resultados de la ciencia microeconómica [\authorfont{Ver Tema 3.A.22}]: la asignación del equilibrio competitivo es óptima en el sentido de Pareto en ausencia de fallos de mercado.

Aunque este resultado tiene un gran valor en sí mismo, depende los siguientes supuestos que pueden resultar muy restrictivos:
\begin{lnum}
    \item \textbf{Principio de precios comunes}. Todos los agentes económicos se enfrentan al mismo sistema de precios y no existen impuestos;
    \item \textbf{Principio de mercados completos}. El bien no numerario dispone de un precio y de un mercado. No existen externalidades ni bienes públicos;
    \item \textbf{Competencia perfecta}. Ningún agente puede afectar de ninguna manera el precio del bien no numerario. Es decir, desde el punto de vista de las decisiones individuales estos son fijos. Ello implica que todos los agentes son precio-aceptantes;
    \item Supuestos sobre preferencias y conjuntos tecnológicos:
    \begin{lnum}
        \item Rendimientos no crecientes a escala, pues en caso contrario no existe solución al problema de maximización de beneficios;
        \item Insaciabilidad local de las preferencias.
    \end{lnum}
\end{lnum}

\subsection{Inobservabilidad de la Funciones Objetivo}
Uno de los supuestos implícitos por los que se cumple es la observabilidad de los agentes de los factores que forman parte de las funciones objetivo de los agentes (p.ej. de los costes de las empresas). Sin esta observabilidad, es decir, con información imperfecta, se podría dar una asignación no eficiente:\footnote{%
    Nótese que la información imperfecta es una condición necesaria pero no suficiente para que se dé una asignación ineficiente.
}
\begin{la}
    \item Si esta información es incompleta pero simétrica.\footnote{%
        Nótese que estaríamos ante un escenario como el estudiado en equilibrios bajo riesgo e incertidumbre [ver Tema 3.A.10], en el que todos los agentes sufren esta particularidad, por lo que no existe un conflicto de intereses sino una voluntad de minimizar un riesgo derivada de las probabilidad de los diferentes estados de naturaleza.
    } según el Teorema de ARROW-HAHN se puede llegar a una asignación óptima si existen mercados contingentes en los que queden estipuladas todas las contingencias;\footnote{%
        Con información simétrica, todos los agentes comparten las mismas creencias sobre los estados de naturaleza; por tanto, el problema se puede representar como un intercambio bajo incertidumbre con agentes ex ante idénticamente informados. En ese marco, el Primer Teorema Fundamental del Bienestar sigue aplicando: los precios reflejan toda la información común y las asignaciones competitivas maximizan la utilidad esperada agregada.
    }
    \item Si este información es incompleta y asimétrica, no existe posibilidad de aseguramiento perfecto frente a una contingencia cuya existencia (selección adversa) o posible realización (riesgo moral) depende de información oculta (selección adversa) o discrecionalidad (riesgo moral) de un agente que:
    \begin{la}
        \item Puede tener incentivos a revelarla (i.e. no existe conflicto de intereses): en cuyo caso los individuos intercambiarán información automáticamente de acuerdo con el principio de revelación de MYERSON y el problema de información asimétrica quedará resuelto;\footnote{%
            De acuerdo con el principio de revelación el individuo solo dice la verdad si es su estrategia óptima, o dicho de otra forma, para cualquier equilibrio bayesiano de Nash en un juego de información incompleta, puede ser representado por un mecanismo directo de incentivos compatibles que tenga un equilibrio en el que los jugadores reporten verazmente sus tipos.
        } o,
        \item Puede no tener incentivos a revelarla (i.e. existe conflicto de intereses) en el caso de que considere que ocultar información le puede traer un beneficio.
    \end{la}
\end{la}

Por lo tanto, en el caso de información asimétrica y conflicto de intereses entramos en el caso de incompletitud de mercados que genera una asignación ineficiente.

\subsection{Economía de la Información: Modelos de Equilibrio Parcial}
La economía de la información examina los efectos de la presencia de información asimétrica y conflicto de intereses: ¿qué sucede cuando los agentes no disponen de toda la información relevante, y algunos de ellos disponen de más información que otros? 
Dos manifestaciones especialmente relevantes de esta suboptimalidad son la selección adversa; y el problema de agencia (y, su manifestación más importante, el riesgo moral) y , tal y como es caracterizada por SALANIÉ (2005), podemos decir que el marco comón en la Econoomía de la Información es:
\begin{lnum}
    \item En su mayoría son modelos de equilibrio parcial, aislando mercados concretos;
    \item Los modelos describen interacciones entre un número reducido de actores, habitualmente solo dos (uno de los cuales suele ser la parte informada);
    \item Los modelos sintetizan las restricciones impuestas por el marco institucional vigente a través de un contrato. Este contrato puede ser explícito (en forma de acuerdo escrito) o implícito (dependiente de un sistema de normas de comportamiento). Un contrato explícito puede ser garantizado por una tercera parte independiente (tribunal, mediador, árbitro…) o por el deseo de las partes de crear y mantener una reputación; y,
    \item Los modelos en origen usan intensivamente la teoría de juegos no cooperativos con información imperfecta, aunque el proceso de negociación habitualmente permite usar un sencillo artilugio, el modelo Principal-Agente, que se explicará a continuación.\footnote{%
        Se trata de juegos bayesianos, en los que los participantes tienen unas creencias a priori sobre la información que desconocen, y revisan racionalmente estas creencias en función de cómo se desarrolla la interacción. Sus soluciones pueden ser equilibrios bayesianos perfectos o equilibrios secuenciales.
    }
\end{lnum}
	
\subsubsection{El Modelo Principal-Agente}
El modelo o paradigma Principal-Agente es la herramienta fundamental de la Teoría de la Agencia. Dicha teoría, iniciada en economía por ROSS (1973), estudia la relación entre dos individuos u organizaciones, el Principal y el Agente, en la cual el Principal trata de diseñar un sistema de compensación (contrato) que incentive al Agente a actuar o decidir en interés del Principal (STIGLITZ, 2008). El problema surge pues cuando existe una asimetría de información entre Principal y Agente y a esto se une el desalineamiento de los intereses de ambos. El objetivo será identificar las características de esta relación, especialmente aquellas relevantes para la distribución de información, a fin de establecer qué tipos de contratos son óptimos para cada situación y qué variables los determinan, siendo la clave el conflicto de intereses, pues los objetivos del Agente se oponen a los del Principal. Por ejemplo, en la relación de agencia entre empleador y trabajador, el salario es un coste para el empresario, pero un ingreso para el trabajador. Y el esfuerzo sería lo opuesto, un coste para el trabajador, pero beneficioso para el empleador.

\paragraph{Taxonomía de los problemas de información asimétrica}
Existen diferentes fuentes o formas de asimetrías informativas. En función de cómo se manifieste esta asimetría, en concreto:
\begin{la}
    \item Según el momento en que se produzca tal asimetría (antes o después de la firma del contrato); y,
    \item Según aquello sobre lo que verse la asimetría, siguiendo la clasificación de ARROW, bien a informaciones o estados de la naturaleza (lo que la parte informada, típicamente el Agente es o sabe: modelos de información oculta), bien referida a acciones/decisiones (lo que la parte informada hace: modelos de acción oculta).\footnote{%
        ARROW, KENNETH (1985), The economics of agency. En Principals and agents: The structure of business, ed. J. Pratt y R. Zeckhauser. Boston: Harvard Business School Press.
    }
\end{la}

En consecuencia, se pueden clasificar las diferentes situaciones en tres grandes categorías:
\begin{la}
    \item Cuando la asimetría se produce ex ante, previa a la firma del contrato, se dan los casos de selección adversa o filtrado  (screening) y señalización (signalling). En ambos casos son problemas de información oculta;\footnote{%
        Aunque esta terminología es común en la literatura especializada, ciertos autores como MAS-COLELL, WINSTON y GREEN (1995) o el análisis seminal de AKERLOF prefieren reservar el término “selección adversa” para referirse al problema o resultado de la asimetría informativa cuando ningún participante busca atajarla activamente, mientras que se refieren a modelos de filtrado y de señalización cuando una de las partes (el Principal y el Agente, respectivamente) actúan para mitigar el problema a través del diseño de un contrato de autoselección o de señales.
    } y,
    \item Cuando la asimetría informativa se manifiesta ex post, después de la firma del contrato, se dan casos de riesgo moral . Pueden ser problemas de información oculta o de acción oculta.\footnote{%
        No obstante, algunos autores optan por distinguir aquí si la asimetría informativa se refiere a una acción o decisión relevante, a lo que llaman riesgo moral en un sentido más restringido, o si se refiere a una información relevante solo conocida por el Agente (Mas-Colell, Winston, \& Green, 1995).
    }
\end{la}




\clearpage
\section{Riesgo moral}
\puntosclave{
    Las asimetrías informativas posteriores a la firma del contrato dan lugar a los modelos de riesgo moral, ya sean de acción oculta o de información oculta a la contraparte
    
	En ambos casos se produce un conflicto entre el reparto eficiente de riesgos y la provisión adecuada de incentivos cuando una parte posee información privada relevante, pues conseguir uno de estos objetivos implicará renunciar al otro
    
	Se alcanzan típicamente soluciones que permiten que la acción o información ocultas sean implícitamente concidas verazmente por la parte no informada del contrato, pero para incentivarlo se debe distorsionar el reparto de riesgos (pérdidas de eficiencia), por lo que típicamente la solución es subóptima
    
	Existen múltiples extensiones que matizan y enriquecen las conclusiones alcanzadas, así como numerosas aplicaciones a mercados o situaciones concretas
}

Un caso de riesgo moral se caracteriza por:
\begin{lnum}
    \item Una asimetría informativa ex post, que se producirá tras la firma del contrato entre Principal y Agente;
    \item Esta información privada se refiere a una acción relevante (acción oculta) o información relevante (información oculta), que la parte informada, típicamente el Agente, conoce con certeza solo tras la firma del contrato. Sin embargo, esta información o acción no es verificable (y por tanto no puede incluirse en el contrato) y como mucho la parte desinformada, típicamente el Principal, la puede conocer solo de manera aproximada a través de una distribución de probabilidades a priori. Esto implica que necesariamente debe existir una heterogeneidad de informaciones o de acciones posibles. No obstante, sí se podrá conocer por todas las partes alguna consecuencia o resultado relacionado con dicha acción o información oculta, aunque solo puede existir una relación imperfecta, típicamente estocástica, entre la acción o información oculta y el resultado observable;\footnote{%
        Si dicha relación fuera determinística, se podría inferir perfectamente la acción o información oculta a partir del resultado, lo que equivaldría a una situación de información simétrica.
    } y,
    \item Esta acción o información privada no puede ser revelada verazmente de manera directa (simplemente anunciándola), pues en función del modelo ciertos participantes sí desearían revelar de manera honesta, mientras que otros tipos tienen incentivos a (se beneficiarían de) declarar una acción o una información falsa. Por ello la revelación veraz supondrá necesariamente un coste, para que sea a prueba de falseamiento.
\end{lnum}

La diferencia esencial del problema de riesgo moral respecto a las situaciones de asimetría ex ante -que veremos después- consiste en incorporar un problema de reparto de riesgo entre Principal y Agente. 

\textbf{NOTA.-} Esto obliga a una breve mención previa al problema general de reparto de riesgos: el reparto de riesgos óptimo en condiciones de simetría informativa tiene su propio desarrollo, pero sus resultados pueden resumirse en que el riesgo debe recaer relativamente más en aquellos agentes más dispuestos a (o menos perjudicados por) soportarlo.\footnote{%
    En caso de un contrato bilateral, si una parte es aversa y el otro no la lógica es simple: si en este caso el averso soportara algo de riesgo, siempre se podría trasladar riesgo del averso a la otra parte del contrato, lo que no le supone perjuicio alguno (si es amante incluso le supone un beneficio). Se alcanza así una mejora en sentido de Pareto (alguien mejora sin empeorar a nadie). Las mejoras de Pareto se agotan y se alcanza el óptimo cuando ya no es posible reasignar riesgo entre las partes, en este caso, cuando el averso ya no soporta nada de riesgo y su contraparte lo soporta todo. En caso de que ambas partes fueran aversas al riesgo, no quedaría más remedio que ambas soportaran parte del riesgo en el reparto óptimo, cada una en función inversa a su aversión al riesgo comparativa (el menos averso debería soportar relativamente más riesgo).
}

El problema de riesgo moral aparece cuando este reparto de riesgos se ve afectado por la asimetría informativa, en sentido de que tras la firma del contrato una de las partes pueda bien realizar acciones ocultas a la contraparte, bien conocer informaciones ocultas a la contraparte, que afecten al resultado de la relación contractual que deba repartirse entre ambas. 

La eficiencia (aquí, reparto de riesgo) se ve comprometida por la provisión de incentivos a que la parte informada revele verazmente la información privada.  Por lo tanto, el Principal diseñará el contrato anticipando esta asimetría para conseguir alinear los intereses de las partes, es decir, incentivando que la acción oculta sea la solicitada o que se revele verazmente la información oculta.\footnote{%
    Puesto que en realidad, lo quese pretende es aumentar la utilidad esperada de ambas partes mediante la introducción de incentivos. Si no, un Agente averso al riesgo estaría satisfecho con un pago mínimo que satisfaciese su nivel de utilidad de reserva.
}

Se desarrollará a continuación el caso de acción oculta. La representación temporal genérica del problema de riesgo moral (siendo la acción oculta el esfuerzo) sería:

\imagenfit{Esquem_Temporal_RM_AccionOculta.png}{Eje temporal estándar de un problema de riesgo moral o acción oculta}{Apuntes ICEX-CECO. Tema 3.A.13}

\subsection{Estructura \textit{baseline}}
Se parte de un Principal que contrata a un Agente, para que lleve a cabo en su nombre una tarea, que a menudo se representa como esfuerzo ($e$), a cambio de una remuneración o contrapartida. Fruto de esta:

\begin{center}
\renewcommand{\arraystretch}{1.25}
\setlength{\tabcolsep}{6pt}
\begin{tabular}{|p{0.18\linewidth}|p{0.39\linewidth}|p{0.39\linewidth}|}
\hline
 & \textbf{Respecto del Agente} & \textbf{Respecto del Principal} \\
\hline
\textbf{¿Qué quieren?} &
Maximizadores de Utilidad -- {\scriptsize Los Agentes son maximizadores de su utilidad. Poseen un nivel de utilidad mínimo o utilidad de reserva, $\bar U$, que será la utilidad mínima que estarán dispuestos a obtener si firman el contrato. Esta utilidad dependerá de dos elementos: (i) un bien, que es la remuneración, $w(x)$; y, (ii) un mal, que es la desutilidad del esfuerzo, $v(e_i)$. Para el modelo baseline, simplifica mucho las cosas si su función de utilidad es separable: $u(w(x)) - v(e_i)$.} &
Maximizadores de Beneficios -- {\scriptsize Entendido como el resultado bruto de la relación menos el pago al Agente, $x - w$. No obstante, el resultado bruto $x$ es incierto y está sujeto a una distribución de probabilidad.} \\
\hline
\textbf{¿Qué son?} &
Aversión al riesgo -- {\scriptsize Debe ser para estar ante un auténtico problema de riesgo moral. Si no fuera así, no habría conflicto entre el reparto de riesgo y provisión de incentivos, ya que se le haría soportar todo el riesgo a fin de incentivar un mayor esfuerzo.} &
Indefinido -- {\scriptsize A diferencia del Agente, su actitud frente al riesgo puede ser de cualquier tipo (aversión, neutralidad o amor por el riesgo).} \\
\hline
\textbf{¿Qué información poseen?} &
Información privada -- {\scriptsize No será observable para el Principal, y con ello tampoco lo será cuánta desutilidad o coste le produce el esfuerzo a cada Agente, $v(e_{i})$. Es decir, cada Agente conoce con certeza su esfuerzo y la desutilidad asociada, pero el Principal no.} &
Distribuciones de probabilidad asociadas al esfuerzo -- {\scriptsize Se tiene que cada nivel de esfuerzo (o acción) ejercida por al Agente afecta a la distribución de probabilidades que rige dicho resultado bruto para el Principal y el Agente, $x$. De este modo se tiene una función de densidad de probabilidad de $x$ condicionada al nivel de esfuerzo: $\mathrm P[x =x_n\;|\;e_i]=p_n(e_i)$ . En principio, no se puede descartar ningún nivel de resultado para un ningún nivel dado de esfuerzo (todos los niveles tienen alguna probabilidad, aunque sea mínima, de alcanzar cualquier resultado).} \\
\hline
\textbf{¿Qué eligen?} &
Grado de esfuerzo -- {\scriptsize Entre una heterogeneidad de acciones posibles o niveles de esfuerzo, que pueden presentar diferentes tipos}\footnote{%
    Los modelos más sencillos son dicotómicos (solo hay dos niveles de esfuerzo, alto o bajo y dos estados de la naturaleza, favorable y desfavorable), lo que los simplifica enormemente e incluso permite su versión gráfica. En general pueden darse más de dos niveles y dos estados, incluso puede tenerse un continuo (conjunto infinito) de niveles de esfuerzo y de estados, acotados entre un nivel mínimo y uno máximo.
} &
El pago al Agente $\mathbb{E}[w]$ -- {\scriptsize Sabemos que este resultado viene en parte condicionado por el esfuerzo del Agente, luego $\mathbb{E}[w(x)\;|\; e_i]$. La razón para esto será, en su caso, incentivar al Agente a aplicar un determinado esfuerzo que aumente las probabilidades de mayores resultados brutos.} \\
\hline
\textbf{¿Cómo eligen?} &
Maximizará su Utilidad Esperada (v.N-M) -- {\scriptsize Se comportará según lo previsto en la teoría de la utilidad esperada, es decir, maximizará su función de utilidad v.N-M, $\mathbb{E}[u(w(x)) - v(e_i)]$} &
Maximizará su Utilidad Esperada (v.N-M) -- {\scriptsize El Principal maximizará su utilidad v.N-M o beneficio neto esperado, $\mathbb{E}[x - w] = \mathbb{E}[x] - \mathbb{E}[w]$.} \\
\hline
\end{tabular}
\end{center}

\subsubsection{Planteamiento del problema de optimización}
El problema de optimización reside en la definición del plan óptimo de incentivos que maximice la función de utilidad esperada del Principal. Por tanto, el Principal determinará el nivel de esfuerzo óptimo que solicitará al Agente (pero que no puede exigir en el contrato por no ser verificable) y el esquema de compensación óptimo para cada estado de la naturaleza, $\{e_i,w(x)\}$, de modo que el Agente elija, en caso de aceptar el contrato, ejercitar el esfuerza que le resulta más beneficioso, dado el plan de incentivos o menú de contratos. La función objetivo del Principal para el caso \textit{baseline}, consistente en maximizar su beneficio neto (de remuneración al Agente) esperado quedaría así:

\eqblock{
\begin{aligned}
    \boxed{\begin{aligned}
        \max_{\{e_i,w(x)\}}\quad \mathbb E[\underbrace{x\;|\;e_i}_{p_n(e_i)}]-\mathbb E[w(x\;|\;e_i)]\\[0.5em]
        \text{s.a.}\quad
        \left\{\begin{aligned}
            &\text{RP:}&&\bar U\leq \mathbb E[u(w(x\;|\;e_i))-v(e_i)]\qquad \forall i\\[0.5em]
            &\text{RCI:}&& e\in\arg \max_{\hat e}\{\mathbb E[u(w(w\;|\;\hat e))-v(\hat e)]\}\\[0.5em]
        \end{aligned}\rigt.
    \end{aligned}}
\end{aligned}
}{Modelo baseline: Problema de optimización y restricciones. Riesgo moral y acción oculta}
 
Como vemos, esta optimización estará sujeta al siguiente sistema de restricciones:
\begin{lnum}
    \item \textbf{Restricción de participación}. El Agente debe obtener una remuneración y ejercer un esfuerzo tales que le permitan alcanzar una utilidad esperada de al menos su utilidad de reserva; y,
    \item \textbf{Restricción de compatibilidad de incentivos}. El Agente ejecutará el nivel de esfuerzo que se le pide, es decir, que maximice dicha utilidad esperada cuando ejecuta el esfuerzo solicitado y no cualquier otro. Es decir, el contrato debe estar diseñado de tal forma que al Agente le convenga voluntariamente elegir el esfuerzo que el principal desea que realice. El Agente hará lo que el Principal quiere, si y solo si hacerlo le reporta más utilidad esperada que cualquier otra opción de esfuerzo. El contrato debe alinear los incentivos del Agente con los del Principal.\footnote{%
        La clave es reconocer que busca alinear los intereses de Principal y Agente, dado el potencial conflicto entre reparto de riesgos y provisión de incentivos. Que el Agente soporte parte o todo el riesgo no es consecuencia de un reparto eficiente del mismo, sino la única manera que posee el Principal de incentivar al Agente, ya que sólo éste puede influir en el resultado de la relación vía su esfuerzo. Y el resultado aleatorio �� es la única variable verificable que tiene el Principal como fuente de información indirecta acerca de la acción oculta (esfuerzo) del Agente a fin de incentivarlo: el Principal estará dispuesto a pagar más remuneración al Agente para aquellos resultados que son estadísticamente más probable que sucedan con un esfuerzo mayor que con un esfuerzo menor, pero independientemente de que el resultado sea mayor o menor.
    }
\end{lnum}

\subsubsection{Solución analítica}
La resolución de este problema es en general altamente compleja debido a que contiene una doble maximización: la RCI, es a su vez un problema de maximización. En los casos dicotómicos esta dificultad suele desaparecer, pero en el caso general presenta un reto técnico. 

En un primer momento se adoptó el denominado enfoque de primer orden (first-order approach), consistente en sustituir esta restricción problema por su condición de primer orden (HOLMSTRÖM, 1979).\footnote{%
    Ver: BENGT HOLMSTRÖM (1979). Moral Hazard and Observability. The Bell Journal of Economics, 10(1), 74–91.
} 

\paragraph{Casos particulares: Análisis de las CPO’s (HOMSTRÖM)}
HOLMSTRÖM demuestra que, bajo ciertas condiciones de regularidad, en particular:
\begin{lnum}
    \item Convexidad del conjunto de niveles de esfuerzo posible;
    \item Diferenciabilidad de las funciones cuyo dominio implica al conjunto de niveles de esfuerzo, esto es la función de pago [w(·)], la función de utilidad condicionada al esfuerzo $[u(· | e)]$ y la función de desutilidad del esfuerzo $[v(·)]$;
    \item Concavidad respecto al esfuerzo de la Función de Utilidad Esperada del Agente. Esto implica que cualquier máximo local (derivada igual a cero) es también un máximo global. Es la condición clave que permite sustituir la restricción global de incentivos por la condición local (CPO);
    \item Regularidad de la función de densidad  (de probabilidad del resultado condicionada al esfuerzo). Esto es, un mayor esfuerzo aumenta la probabilidad relativa de resultados altos. Esto garantiza que el incentivo induce efectivamente más esfuerzo;\footnote{%
        Aunque pueda parecer contraintuitivo, a priori nada obliga a que los resultados mejores (mayor $x$) se produzcan con mayor probabilidad con esfuerzos mayores, por lo que la remuneración no tiene por qué ser creciente con el resultado (esto es, no tiene por qué darse $w'(x)>0$). Por otro lado, también sería posible que un mayor esfuerzo elevara la probabilidad de obtener mejores resultados, pero solo hasta cierto nivel de resultados, más allá del cual un mayor esfuerzo apenas mejore la realización de los resultados aún más altos. En este caso, esos niveles intermedios de resultados generarán mayores remuneraciones que los resultados más altos, pues aquellos son más sensibles o dependientes del esfuerzo (MAS-COLELL, WINSTON, y GREEN, 1995).
    } y,
    \item No existencia de \textit{corner solutions} (esfuerzos interiores). El esfuerzo óptimo se encuentra en el interior del dominio de esfuezos posibles, no en la frontera de éste. Derivada primera será igual a cero y derivada segunda negativa.
\end{lnum}

En lugar de imponer toda la restricción de incentivos (que es global y complicada), podremos imponer únicamente la condición de primer orden (CPO’s) del problema del Agente.

Esto es, incorporaremos la CPO del problema de maximización de incentivos al problema de optimización del Principal asumiendo que el Agente, en el óptimo, no tendrá incentivos para desviarse, por lo que, estaremos ante la situación siguiente:

$$\begin{aligned}
    \frac{\partial}{\partial e}\mathbb E[u(w(x\;|\;e))-v(e)]=0\\[0.5em]
    \mathbb E[u'(w(x\;|\;e))\frac{\partial w(x\;|\;e)}{\partial e}+u'(w(x\;|\;e))\frac{\partial x}{\partial e}]=v'(e)  
\end{aligned}$$

Pero, como podemos asumir que el resultado, o más concretamente, la ganancia en función del resultado depende positivamente del esfuerzo empleado, podemos reescribir la CPO así:

$$\begin{aligned}
    \mathbb E[u'(w(x))x_e]=v'(e)\qquad \text{donde,}\quad x_e=\frac{\partial x}{\partial e}
\end{aligned}$$

Es decir, el Pago Marginal esperada del esfuerzo será igual a la Desutilidad Marginal de dicho esfuerzo. Por lo tanto, el Problema de Optimización del Principal se formulará de la siguiente manera:

$$\mathcal L=\mathbb E[x-w(x)]+\mu\mathbb (E[u(w(x\;|\;e))-v(e)]-\bar U)+v(\mathbb E[u'(w(x\;|\;e))x_e]-v'(e))$$
 
Donde $\mu$ será el Multiplicador asociado a la Restricción de Participación y $v$ será el Multiplicador asociado a la Restricción de Compatibilidad de Incentivos. 

\paragraph{Solución general: Método dos etapas (GROSSMAN y HART)}
Para todos los demás casos es necesario recurrir al proceso en dos etapas de GROSSMAN y HART (1983), que permite analizar las características del contrato óptimo independientemente de si se es capaz de identificar el nivel de esfuerzo óptimo o no:
\begin{lnum}
    \item Primero, para cada nivel de esfuerzo que se desee contratar, el Principal debe encontrar un esquema de remuneración $w^{*}_{x}$ que incentive la participación y la compatibilidad de incentivos (esfuerzo requerido) para cada Agente. Es decir, debe determinar cuánto le va a costar contratar al Agente y cuánto le costará solicitar cada nivel de esfuerzo;\footnote{%
        ¿Qué factores estructurales determinarían cómo de costoso es incentivar cada esfuerzo u acción oculta? El coste de incentivar aumentaría con el grado de aversión al riego del Agente (lo costoso de hacerle soportar más riesgo), con la desutilidad relativa de mayores niveles de esfuerzo o con lo débil que sea la relación entre esfuerzo y resultado (esto es, cuando el resultado sea poco sensible el esfuerzo).
    } y,
    \item Conocido dicho esquema de remuneración, el Principal debe elegir el nivel de esfuerzo (o de acción oculta, en general) que maximice su beneficio esperado, $e^*$.
\end{lnum}

\subsubsection{Propiedades notables}

\paragraph{La no-linearidad del esquema de compensación: Compatibilidad de Incentivos}
Respecto a la RCI cabe destacar que dejar que el Agente soporte riesgo no busca que se esfuerce más para lograr mejores resultados, sino que es la única herramienta que tiene el Principal para inferir e incentivar un nivel de esfuerzo concreto. En conclusión, el esquema de compensación óptimo $w^{*}_{x}$ no adoptará típicamente una forma sencilla (v.g. la lineal), pues será una función del contenido informativo de los diferentes niveles de resultado $x$ (es decir, de cuán probable es inferir un nivel de esfuerzo dado un resultado observado), lo que es improbable que varíe con $x$ de una manera sencilla en la mayoría de los problemas (MAS-COLELL, WINSTON, y GREEN, 1995).

\paragraph{Restricción de participación: Saturación}
Respecto a la Restricción de Participación cabe destacar que:
\begin{lnum}
    \item En un modelo como el presentado esta se satura, reflejando que el Agente, tanto si realiza niveles de esfuerzo altos o bajos, obtendrá en todo caso su utilidad de reserva en términos esperados, lo cual es consecuencia en este caso de su nulo poder negociador, siendo el Principal quien se apropie y maximice el excedente;\footnote{%
        No obstante, en casos muy poco habituales esta restricción no tiene por qué saturarse (por ejemplo, cuando la utilidad del Agente no es separable en el salario y el esfuerzo) (LAFFONT y MARTIMORT, 2002).
    }
    \item En el contrato óptimo el Agente debe soportar ineficientemente un mayor riesgo. Esto supone que, dada una utilidad de reserva que debe ser alcanzada por el Agente en términos esperados, un mayor riesgo implicará que el Agente deba ser compensado por ello con una mayor remuneración. Es decir, si bien el Agente sigue obteniendo su utilidad de reserva de media, ahora ello solo puede ser obtenido con un pago extra o compensación por la desutilidad de un mayor riesgo asumido. En consecuencia, la parte informada en este modelo no recibirá algo así como una renta de información derivada de su ventaja informativa, simplemente un pago extra o compensación por un mayor riesgo asumido, pero con igual utilidad esperada que en el caso de información simétrica. 
\end{lnum}

\paragraph{La actitud frente al riesgo del Agente: Tipos de esquemas compensatorios}
La actitud frente al riesgo del Agente juega un papel fundamental, puesto que:
\begin{la}
    \item En caso de que el Agente fuera neutral al riesgo, el incentivo al esfuerzo puede obtenerse sin comprometer el reparto eficiente de riesgos, pues la solución óptima es que el Agente soportara todo el riesgo a modo de “franquicia” (mientras el Principal obtendría un pago fijo independiente del estado de la naturaleza). No existiría un problema de riesgo moral, pues este reparto de riesgos también es eficiente y sería el mismo que se daría bajo información simétrica;
    \item En caso de que fuera averso al riesgo sí habría un problema de riesgo moral, pues el Agente soportará más riesgo de lo que sería eficiente 45, pero solo para conseguir que se esfuerce lo requerido. Solamente en el caso de que el esfuerzo óptimo requerido por el Principal fuese el más bajo posible podría lograrse un reparto eficiente de riesgos, pues basta la mínima remuneración posible en todos los estados de la naturaleza (aseguramiento perfecto) para asegurarse dicho esfuerzo mínimo. La distorsión o pérdida de eficiencia se produce cuando se quiere inducir un nivel de esfuerzo  estrictamente superior al mínimo posible.
\end{la}

\paragraph{Análisis del Bienestar}
Desde el punto de vista del bienestar, el contrato óptimo del problema de riesgo moral también supone un \textit{second best}, pues si bien típicamente el Principal conseguirá inducir con certeza un nivel de esfuerzo en el Agente, tal como si no hubiera asimetría informativa, esto se conseguirá mediante una pérdida de eficiencia. El Agente tendrá garantizado el nivel de utilidad de reserva, igual que en el caso de información simétrica, por lo que no sufre pérdidas de bienestar (si bien para lograrlo requiere ahora una mayor remuneración que compense el mayor riesgo soportado). Sin embargo, el Principal tendrá que distorsionar el reparto de riesgos si quiere incentivar al Agente a realizar esfuerzos estrictamente superiores al mínimo posible, de modo que deba dedicar parte del excedente a compensar al Agente por el mayor riesgo asumido bajo el contrato óptimo.\footnote{%
    Con la excepción del caso en que el nivel de esfuerzo óptimo para el Principal sea el mínimo posible, como se ha visto
}


\subsection{Ejemplo: Desarrollo gráfico}
Supongamos un problema de reisgo moral donde existen dos resultados posibles $x_{1}$ y $x_{2}$ tal que $x_{2} > x_{1}$. La probabilidad de dichos resultados viene dada por una variable aleatoria de estado de naturaleza sobre la cual afecta el esfuerzo o desempeño que ejerce el Agente, que puede escoger entre un esfuerzo elevado ($e^{H}$ o un esfuerzo bajo ($e^{L}$)

$$\begin{aligned}
    &\mathrm P\left[x=x_2\;|\;e^H\right]=p^H &&\wedge  &&\mathrm P\left[x=x_1\;|\;e^H\right]=1-p^H\\[0.5em]\\
    \hline\\
    &\mathrm P\left[x=x_2\;|\;e^L\right]=p^L &&\wedge &&\mathrm P\left[x=x_1\;|\;e^L\right]=1-p^L\\[3em]
    &&&\hspace{-1.6em}p^H>p^L
\end{aligned}$$
 
Realizar un esfuerzo elevado o un esfuerzo bajo también influye sobre el Agente, ya que este percibe una desutilidad de este tal que $v(e^{H}) > v(e^{L})$:

$$\begin{aligned}
    &\max_{\{w_1,w_2,e\}}\quad \overset{\substack{\text{\scriptsize Expexted Utility (discreta)}}\\[1em]}{p^H(x_2-w_2)+(1-p^H)(x_1-w_1)}\\[1em]
    &\text{s.a.}\quad
    \begin{aligned}
        &\text{R.P:} &&0\leq p^H+(1-p^H)u(w_1)-v(e^H)\\[0.5em]
        &\text{RCI:} &&\underbrace{p^Lu(w_2)+(1-p^L)u(w_1)-v(e^L)}_{\substack{\text{\scriptsize Expected Utility con $e^H$}}}\leq \underbrace{p^Hu(w_2)+(1-p^H)u(w_1)-v(e^H)}_{\substack{\text{\scriptsize Expected Utility con $e^L$}}}
    \end{aligned}
\end{aligned}$$

De esta forma, para que $f(w_1, 2_2)$ represente el conjunto de contratos ($w_2,w_1$) tal que la restricción de compatibilidad de incentivos se cumple:

$$\begin{aligned}
        \begin{aligned}
            &\underset{\substack{\\[1em]\\(\checkmark:\;\text{RCI})}}{(w_2,w_1)\in f(w_2,w_1)}&&\iff &&p^Hu(w_2)+(1-p^H)u(w_1)-v(e^H)=p^Lu(w_2)+(1-p^L)u(w_1)-v(e^L)
        \end{aligned}\\[2em]
\end{aligned}$$

Debemos destacar dos cosas:
\begin{lnum}
    \item Para todo conjunto de contratos que pertenezca a la función que hace compatible la RCI, el agente será indiferente entre llevar a cabo un esfuerzo alto o un esfuerzo bajo: $\forall (w_2,w_1)\in f(w_2,w_1)$, $e^H\sim e^L$; y,
    \item Si dado un conjunto de contratos que cumple con la condición, aumentamos alguno de sus componentes, $\uparrow\; w_2$ ($\uparrow\;w_1$), la igualdad anterior se convierte en una desigualdad $>$ ($<$)
\end{lnum}

\paragraph{Menú de contratos y RCI: Agente}
Podremos graficar estas condiciones de manera que, para el Agente:

\imagenfit{RiesgoMoral_EjGrafico.png}{Restricción e compatibilidad de incentivos ($f(w_2,w_1)$ y curvas de indiferencia del agente)}{Elaboración propia}

Por lo tanto, como vemos:
\begin{la}
    \item Tras un incremento en el primer contrato, $\uparrow w_1$, nos situaríamos en el punto A, donde el Agente tendrá más utilidad dedicando un esfuerzo bajo al cometido ($e^L$);
    \item Por el contrato, si incrementamos el segundo contrato, desde la curva de compatibilidad, $\uparrow \;w_2$, el Agente obtendrá más utilidad dedicando un mayor esfuerzo al acometido ($e^H$).
\end{la}

Por un razonamiento análogo podemos llegar a la conclusión general: (i) en la región superior, $\{w_2,w_1\}>f(w_2,w_1)$, el Agente obtiene más utilidad dedicando menos esfuerzo; y, (ii) en la región inferior, $\{w_2,w_1\}<f(w_2,w_1)$, el Agente obtiene más utilidad dedicando un mayor esfuerzo.

\paragraph{Beneficio: Principal}
Por lo que respecta al Principal, el beneficio que obtendrá menos el salario que deberá pagar por el resultado también se ve afectado por el esfuerzo, por lo que deseará incentivar un mayor esfuerzo por parte del Agente:

\imagenfit{Beneficio_EjRiesgoMoral_Principal.png}{Beneficio obtenido por el principal en función del esfuerzo del agente y el salario pagado}{Elaboración propia}

La zona sobre la que su beneficio esperado es el mismo independientemente del esfuerzo del Agente es:

$$p^H(x_2-w_2)+(1-p^H)(x_1-w_1)=p^L(x_2-w_2)+(1-p^L)(x_1-w_1)$$

Por lo tanto, por el gráfico y la bisectriz de indiferencia, sabemos que:
\begin{la}
    \item Si aumentamos el diferencial ($x_1-w_1$), nos desplazaremos hacia el punto A. En términos económicos esto representa el surplus que recibe el Agente dado un bajo nivel de esfuerzo; y,
    \item Si aumentamos el diferencial ($x_2-w_2$), nos desplazamos hacia el punto B. En términos económicos, esto respesnta el surplus que recibe el agente dado un esfuerzo alto.
\end{la}

De forma análoga para las regiones (en función del lado de la bisectriz). Por tanto, el Principal preferirá todos los menús de contratos donde: $\forall \{(x_2-w_2), (x_1,w_1)\}\;:\;(x_2-w_2)>(x_1-w_1)$.

\paragraph{Menú de contratos óptimo: Caja de EDGEWORTH}
Si sitúamos los dos gráficos superpuestos apoyándonos en la Caja de EDGEWORTH, tenemos:

\imagenfit{Optimo_Conjunto_RM_EjGrafico.png}{Caja de EDGEWORTH: Menú de contratos óptimos para el Principal}{Elaboración propia}

Si el principal decide ofrecer las condiciones de $L$, entonces el Agente escogerá realizar un esfuerzo bajo ($e^L$), ya que $L>f(w_2,w_1)$. De la misma forma, si el principal decide escoger el contrato $H$, el Agente realizará aún así un esfuerzo bajo ($e^L$) alcanzando un mayor nivel de utilidad. Entonces, ¿qué debe escoger el Principal? Como vemos en el gráfico de la izquierda, el contrato óptimo para el Principal se encuentra en el punto de intersección entre las curvas de indiferencia, $H'\in f(w_2,w_1)$. De hecho, si nos fijamos en la bisectriz del plano de beneficios del principal, este contrato $H'$ queda por debajo de la bisectriz por lo que, independientemente de lo que haga el Agente (que en este punto es indiferente ante una esfuerzo alto o bajo), se trata del mejor contrato calidad precio y aún dispondría de margen de mejora para incentivar un mayor esfuerzo en el Agente.

Ahora ya sabemos cuál es el menú de contratos que debe ofrecer el Principal para cada nivel de esfuerzo: $L(e^L),\;H'(e^H)$. 

\subsection{Aplicaciones en diversos contextos económicos}
Como ejemplos de aplicaciones de problema de riesgo moral en mercados o situaciones concretas, se puede citar:
\begin{lnum}
    \item La regulación de empresas también puede verse afectada por situaciones de riesgo moral, por ejemplo, cuando se busca que la empresa regulada busque mejoras (reducciones) en costes. Con la regulación del coste de servicio basada en costes históricos se corre el riesgo de que la empresa pierda todo interés en reducir sus costes, pues su cobertura está asegurada, así que puede permitirse que la empresa se beneficie al menos parcialmente de la reducción de costes que consiga alcanzar;
    \item La relación laboral entre trabajador y empresa (SHAPIRO y STIGLITZ, 1984), en la que la empresa debe asegurarse de que el trabajador se esfuerce, para lo cual puede llevarse a cabo actuaciones de control limitadas. El castigo de ser despedido si el trabajador es descubierto holgazaneando se ve disminuido si el trabajador puede encontrar un nuevo trabajo con facilidad, por lo que el paro actúa como mecanismo disciplinador.\footnote{%
        SHAPIRO y STIGLITZ (1984). Equilibrium Unemployment as a Worker Discipline Device. The American Economic Review. 74 (3): 433–444.
    }
\end{lnum}





\clearpage
\section{Selección adversa}
\puntosclave{
    En función de quién tome la inciativa para mitigar la asimetría de información existente antes de la firma del contrato se tienen los modelos de selección adversa (Principal) y los modelos de señalización (Agente)
    
	En los modelos de selección adversa o filtrado el Principal maximiza su beneficio sujeto a restricciones de participación y de compatibilidad de incentivos del Agente, mediante el diseño de contratos óptimos de autoselección para que el Agente escoja aquél pensado para su tipo y no otro
    
	Pero estos contratos óptimos, si bien consiguen revelar técitamente la información privada de los Agentes, implican una pérdida de eficiencia: existe un conflicto entre la consecución de la eficiencia y la minimización de rentas de información que obtienen ciertos Agentes. Es decir, logra una asignación subóptima (second best)
    
	En los modelos de señalización es el Agente quien emite una señal costosa para revelar verazmente su información privada (hay incentivos a falsearla)
    
	Estos juegos de señalización admiten formalmente equilibrios agrupadores (no se consigue diferenciar entre Agentes) y separadores (se consigue diferenciar entre Agentes), pero en ambos casos la solución será típicamente subóptima
    
	En ambos casos existen múltiples extensiones que matizan y enriquecen las conclusiones alcanzadas, así como numerosas aplicaciones a mercados o situaciones concretas
}
	
Dentro de la categoría de modelos con asimetrías ex ante (previas a la firma del contrato) se pueden distinguir a su vez los modelos de selección adversa o filtrado (screening), en los que es la parte desinformada (Principal) quien trata de mitigar dicha asimetría, y los modelos de señalización (signalling) en los que es la parte informada (Agente) quien busca un mecanismo para revelar verazmente su información privada.
Un caso de selección adversa se caracteriza por:
\begin{lnum}
    \item Una asimetría informativa ex ante, que precede a la firma del contrato entre Principal y Agente; 
    \item Esta información privada o tipo se refiere a una característica o dato relevante sobre lo que una de las partes de la relación es o sabe (información oculta). Sin embargo, esta información no es verificable, a lo sumo la contraparte la conoce solo de manera aproximada: como distribución de probabilidades (parte desinformada);\footnote{%
        Esto permitiría utilizar juegos bayesianos para llegar a un equilibrio parcial: Equilibrio Baye-Nash. Sin embargo, supone un supuesto demasiado exigente puesto que se parte de la base de que la ditstribución de probabilidades de tipos es “common-knowledge”; esto es, fundamento matemático acuñado por ROBERT AUMANN : es conocimiento común cuando todos los agentes lo saben, todos saben que los demás lo saben, todos saben que todos saben que los demás lo saben, y así sucesivamente sin fin. AUMANN mostró que este concepto requiere una clausura infinita del conocimiento mutuo, lo que lo convierte en una base formal para el razonamiento coordinado y las expectativas compartidas en teoría de juegos y economía.
    } y,
    \item Esta información privada no puede ser revelada verazmente de manera directa (simplemente anunciada por los Agentes), pues en función del modelo ciertos tipos sí desearían revelar su tipo, mientras que otros tipos tienen incentivos (se benefician de) a mentir y hacerse pasar por un tipo diferente al que realmente se les ha asignado. Por ello la revelación veraz supondrá necesariamente un coste. La razón: tiene que ser creíble, a prueba de mentiras.
\end{lnum}

El esquema temporal de la relación se define como:

\imagenfit{EsquemaTemporal_SA.png}{Esquema temporal de un problema de Selección Adversa (o autoselección)}{Apuntes ICEX-CECO. Tema 3.A.13}
 
\textbf{NOTA.-} Cabe destacar que, miestras que los modelos de riesgo moral se suelen analizar en el marco de la Teoría de Agencia, no existe semejante uniformidad metodológica para analizar el fenómeno de la selección adversa. La literatura ha dado un tratamiento distinto a los problemas de se selección adversa en función del mercado estudiado.

\subsection{AKERLOFF: mercado de coches de segunda mano}
El modelo seminal de AKERLOF (1970) de coches de segunda mano ilustra las posibles consecuencias de esta asimetría en ausencia de cualquier mecanismo o contrato que busque mitigarla. Aquí los vendedores (Agente) conocen el tipo de coche que ponen a la venta, pero los compradores no (Principal). Si la información fuera simétrica, cada coche se vendería por un precio ajustado a su calidad (resultado eficiente). Pero si existe asimetría de información y solo el vendedor conoce la verdadera calidad de su coche, los compradores, maximizando su utilidad esperada, solo estarían dispuestos a pagar el precio correspondiente a la esperanza de calidad. Esto favorecería a los vendedores de coches de calidad inferior a la media (lemons), que venderían su coche por una cantidad superior a la que les correspondería, mientras que perjudicaría a los vendedores de coches de calidad por encima de la media (peaches). A esto se debe añadir que cada vendedor será consciente de este precio medio que recibirá como máximo, por lo que habrá vendedores de coches de alta calidad que no están dispuestos a venderlos tan baratos (dicho en términos técnicos, su utilidad de reserva en caso de no vender su coche de calidad puede ser superior a ese precio medio). Estos vendedores abandonarán el mercado e incluso, en un caso extremo, solo quedarían a la venta en el mercado los coches de peor calidad (lemons).

Sin embargo, en estos modelos de selección adversa o filtrado, el Principal, conocedor de esta situación, usará su capacidad para diseñar el contrato con el fin de discriminar a los tipos de Agentes, con el fin de que estos encuentren beneficioso revelar su tipo verazmente de manera implícita, a través de sus decisiones. El modelo que se expondrá a continuación incorporará este comportamiento optimizador por parte del Principal.

[IMPORTANTE – REVISAR, AKERLOF NO PLANTEA ESTO QUE VIENE A CONTINUACIÓN]

\subsection{Modelos \textit{screening}: Auto selección}
Se parte de un Principal que contrata a un Agente para que lleve a cabo en su nombre una acción o tarea a cambio de una remuneración o contrapartida.\footnote{%
    En estos modelos de asimetrías informativas ex ante la actitud frente al riesgo juega un papel muy secundario, a diferencia de los modelos de asimetrías ex post (riesgo moral), pues no se produce aquí un problema de reparto de riesgos entre Principal y Agente.
} Esta tarea genera un resultado ($x$) beneficioso para el Principal y servirá para remunerar al Agente.\footnote{%
    También pueden darse muchos casos de selección adversa en los que no hay una acción o cometido explícito del Principal al Agente, sino que la asimetría surge simplemente alrededor de alguna característica innata de este último. Por ejemplo, en el caso del mercado de seguros, el Principal (compañía aseguradora) no le encarga nada al Agente (cliente asegurado), sino que la asimetría se debe a las diferentes probabilidades de sufrir la contingencia asegurada en función del tipo de Agente del que se trate. O en la relación de agencia que se da entre un monopolista (Principal) y sus clientes (Agentes), cuya disposición a pagar aquél desearía conocer a fin de discriminar precios.
} Se incluirá aquí la variable esfuerzo ($e$), siguiendo a MACHO-STADLER y PÉREZ-CASTRILLO, a los solos efectos de tener un marco común y facilitar la comparación con el resto de los modelos del tema:

\begin{center}
\renewcommand{\arraystretch}{1.25}
\setlength{\tabcolsep}{6pt}
\begin{tabular}{|p{0.18\linewidth}|p{0.39\linewidth}|p{0.39\linewidth}|}
\hline
 & \textbf{Respecto del Agente} & \textbf{Respecto del Principal} \\
\hline
\textbf{¿Cómo son?} &
Heterogéneos y múltiples -- {\scriptsize Existe heterogeneidad y pluralidad de Agentes (cada uno con su tipo, denotado por $i$). Los modelos más sencillos son dicotómicos (solo hay dos tipos de Agentes) y podremos ordenar jerárquicamente a los Agentes en función de los más o menos deseables respecto al Principal.} &
Único o múltiples -- {\scriptsize Un elemento clave de estos modelos es si hay un único Principal (\textit{monopolistic screening}) o si hay varios Principales que compiten entre sí por contratar Agentes (\textit{competitive screening}). En el primer caso el Principal retiene todo el poder de negociación, por lo que la remuneración o contrapartida que desembolsará será la mínima indispensable, consiguiendo retener cierto beneficio, mientras que en la presión competitiva el beneficio de los Principales será nulo.} \\
\hline
\textbf{¿Qué quieren?} &
Maximizadores de Utilidad -- {\scriptsize Los Agentes son maximizadores de su utilidad y poseen una utilidad de reserva, $\bar U$. La utilidad obtenida en caso de la firma del contrato viene definida por: (i) la remuneración ($w$) que obtienen y que va en función del tipo que sean, $w_i$; y, (ii) la desutilidad, considerada un mal, como función del esfuerzo ($e$) que variará en función del tipo de agente que sean, $v_i(e)$. Para el modelo base, simplifica mucho las cosas si su función de utilidad es separable: $u(w_i) - v_i(e)$.} &
Maximizadores de Beneficios -- {\scriptsize El Principal será maximizador del beneficio (esperado) y éste vendrá dado por la diferencia entre el resultado obtenido y el pago realizado al Agente: $\mathbb E[\Pi(x(e_i) - w_i]$. Cada nivel de esfuerzo ejercido por al Agente está asociado con un pago o resultado para el Principal, $x$.} \\
\hline
\textbf{¿Qué información poseen?} &
Información privada -- {\scriptsize Si bien el esfuerzo (o cualquier acción) sí será observable en estos modelos para el Principal, no lo será cuánta desutilidad o coste le produce el esfuerzo a cada Agente según su tipo, $v_i(e)$, o cualquier otra característica relevante para el contrato. Es decir, cada Agente conoce su tipo, pero el Principal no.} &
Distribuciones de probabilidad -- {\scriptsize El Principal conoce o intuye la distribución poblacional de los tipos. Además, sabe que el resultado es función del esfuerzo que se emplee en el trabajo, por tanto, podrá expresar el beneficio esperado del negocio como: $\Pi(e)=\sum_{i=1}^Np_ix(e(i))$.} \\
\hline
\textbf{¿Qué eligen?} &
El contrato -- {\scriptsize Los agentes eligen el contrato que más les conviene entre un menú de contratos ofrecido por el Principal.} &
El pago por cada tipo de Agente $w_i$ -- {\scriptsize Típicamente, un mayor esfuerzo (o mejor acción) incrementa las probabilidades de mayores o mejores resultados. Su beneficio esperado vendrá dado por el resultado esperado neto del pago a/los Agentes: $\mathbb E[\Pi(x(e(i)) - w_i]$.} \\
\hline
\textbf{¿Cómo eligen?} &
Maximizará su Utilidad Esperada (v.N-M) -- {\scriptsize Se comportará según lo previsto en la teoría de la utilidad esperada, es decir, maximizará su función de utilidad v.N-M o utilidad esperada, $\mathbb E[u(w_i) - v_i(e)]$.} &
Maximizará su Utilidad Esperada (v.N-M) -- {\scriptsize El Principal maximizará su utilidad v.N-M o beneficio neto esperado, $\mathbb E[\Pi(x(e(i)) - w_i)]$.} \\
\hline
\end{tabular}
\end{center}

\subsubsection{Autoselección: \textit{Monopolistic screening}}
Esta versión del problema se plantea desde la óptica del Principal, quien busca maximizar su beneficio esperado, para lo cual debe:\footnote{%
    Este beneficio es esperado porque el Principal, una vez que identifique a los distintos tipos de Agente, solo cuenta a priori con una distribución de probabilidades sobre los resultados.
}
\begin{lnum}
    \item Diseñar un esquema de remuneración óptimo para cada tipo de Agente ($w_i$); y,
    \item Determinar unos niveles de esfuerzo óptimos exigibles a cada tipo de Agente, $e(i)$
\end{lnum}

\paragraph{Planteamiento del problema de optimización}
¿En qué consiste un contrato óptimo para el Principal? 

El Principal ofrecerá un menú de contratos (en el modelo base, pares de esfuerzo exigido y remuneración para cada tipo, ($e(i), w_i$), de modo que cada tipo de Agente elija libremente el diseñado para él y no otro, es decir, debe ser un mecanismo de auto-selección. La función objetivo del Principal para el caso base quedaría así:

\eqblock{
\begin{aligned}
    \boxed{\begin{aligned}
        \max_{\{e(i),w_i\}}\quad \mathbb E[\Pi(x(e(i))-w_i)]\\[0.5em]
        \text{s.a.}\quad
        \begin{aligned}
            &\text{RP:}&&\bar U\leq u(w_i)-v_i(e[i])\\[1em]
            &\text{RCI:}&& \underset{\forall i,k\in N,\;\;i\neq k}{u(w_k)-v_i(e_k)\leq u(w_i)-v_i(e_i)}\\[1em]
        \end{aligned}
    \end{aligned}}
\end{aligned}
}{Problema de \textit{screening}: Monopolistic screening. Economís de la información: Selección adversa}

A priori podría pensarse que este esquema de contratos es demasiado simple y podría sofisticarse para aumentar el beneficio, bien ofreciendo más contratos que tipos de Agentes (que puedan elegir entre un número mayor de opciones), bien incluyendo algún sistema de comunicación entre Principal y Agente para que este transmita información de la que dependa la remuneración. Sin embargo, nada de esto conseguiría una mejora, conforme al Principio de Revelación, que demuestra que no hay pérdida de generalidad en restringirse a este simple menú de contratos, denominado mecanismo de revelación o directo.\footnote{%
    Principio de revelación: “En cualquier problema de incentivos (como selección adversa), para todo mecanismo que induce cierto resultado de equilibrio cuando los agentes pueden mentir sobre su tipo, existe un mecanismo equivalente (en términos de resultados y pagos esperados) en el que los agentes dicen la verdad sobre su tipo y no tienen incentivos a mentir”. Esto es, para todo mecanismo indirecto que alcance el equilibrio, existe un mecanismo directo (revelación) que alcance ese mismo equilibrio (es equivalente)
} Para ello, estos contratos de filtrado deben cumplir dos condiciones (las restricciones del problema):
\begin{lnum}
    \item \textbf{Restricción de participación}. El Agente al que va dirigido debe obtener, al menos, su utilidad de reserva;
    \item \textbf{Restricción de compatibilidad de incentivos}. El Agente al que va dirigido un contrato en particular ($i$) debe tener incentivos a elegir ese contrato y no cualquier otro ($k$), es decir, debe alcanzar mayor utilidad eligiendo el contrato diseñado para su tipo $i$ que haciéndose pasar por otro tipo y escogiendo un contrato pensado para otro tipo de Agente $k$.
\end{lnum}
 
\paragraph{Condición: Preferencias de los Agentes}
Este problema tiene una serie de características especiales que revelan ciertos rasgos del resultado o solución. Se suele exigir que las preferencias de los Agentes satisfagan la condición SPENCER-MIRRLESS:
\begin{lnum}
    \item \textbf{Condición de signo constante}. La derivada de la relación marginal de sustitución entre esfuerzo y remuneración respecto del tipo debería ser siempre del mismo signo, positivo o negativo; y,
    \item \textbf{Condición de cruce único}. Si se representaran en el plano ($w,e$), dos curvas de indiferencia de dos tipos distintos solo deberían cruzarse una sola vez.
\end{lnum}

Como se puede deducir, estas condiciones son equivalente y su interpretación es intuitiva: en este caso implica que, conforme aumenta el coste o desutilidad del esfuerzo entre Agentes (cuanto más ineficientes), mayor será el salario que deberían recibir para inducirles a un cierto esfuerzo.\footnote{%
    De manera general: un tipo más eficiente también es más eficiente en el margen, de modo que un Agente es más eficiente que otro para cualquier nivel de esfuerzo.
}

\paragraph{Solución: Descomposición en dos etapas y Características}
El problema puede descomponerse en dos etapas, una para cada variable de decisión del Principal.
\begin{lnum}
    \item Determinación de los niveles implementables. Primero, para cada nivel de esfuerzo que se desee contratar, el Principal debe encontrar un esquema de remuneración $w^{*}_i$ que incentive la participación y la compatibilidad de incentivos de cada Agente. Es decir, debe determinar cuánto le va a costar contratar a cada tipo de Agente y qué esfuerzo puede pedirle. Los niveles de esfuerzo para los que exista este esquema se denominan implementables; y,
    \item Una vez encontrados los niveles de esfuerzo implementables, el Principal debe elegir aquél que maximice su beneficio, $e_i^*$. Es muy importante destacar aquí que ni siquiera en los modelos dicotómicos más sencillos tiene por qué ser óptimo para el Principal que todos los Agentes participen. En efecto, inducir tanto la participación como la compatibilidad de incentivos es costoso para el Principal por lo que puede resultarle prohibitivo contratar todo tipo de Agentes, prefiriendo que solo participe una parte de los Agentes (habitualmente los “mejores” o menos costosos).
\end{lnum}

\paragraph{Restricción de Participación: Insaturación} 
No todas se saturan o son vinculantes en un conjunto de contratos óptimo. En este caso concreto, la RP sólo resulta operativa para el Agente con la mayor desutilidad del esfuerzo (el más costoso o peor), por lo que el resto de los Agentes alcanzarán una utilidad superior a la de reserva. Este exceso se explica a continuación.

\paragraph{Restricción de Compatibilidad de Incentivos:}
De su manipulación puede derivarse un resultado importante: para conseguir la compatibilidad de incentivos (que cada Agente elija su contrato), cada tipo de Agente ($i$) recibirá típicamente una remuneración por encima de su utilidad de reserva (salvo el peor en nuestro caso concreto) denominada renta de información. Esta renta representa el exceso de remuneración que un Principal tiene que desembolsar debido a la presencia de asimetría informativa . Su cuantía sería la máxima ganancia que un tipo de Agente obtendría por hacerse pasar por otro tipo distinto. 

En el caso contemplado, para cada tipo de Agente esta renta típicamente se incrementa conforme es más deseable o “mejor” (menor coste de esfuerzo), por lo que la asimetría supone aquí mayores rentas para los “mejores”. Sin embargo, en otros modelos se da el fenómeno contrario: son los agentes menos deseables para el Principal los que se benefician de rentas de información . El resultado del problema es el conjunto de contratos óptimo (desde la óptica del Principal), que consigue que los diferentes tipos de Agentes, que participan en la relación y escogen el contrato para ellos diseñado, se delaten o revelen tácitamente su tipo o información privada. En el óptimo, los Agentes más eficientes o con menos desutilidad del esfuerzo realizarán mayores niveles de esfuerzo. 

\paragraph{Análisis del Bienestar}
En términos de bienestar, es evidente que este sistema de filtrado o auto-selección supone una mejora de eficiencia respecto a una situación en la que no pueda usarse este mecanismo, como en el análisis seminal de AKERLOF (1970). 

Ahora bien, estos contratos de filtrado típicamente mitigan el problema de información, pero sin llegar a recuperar la asignación que se daría en ausencia de asimetría informativa o first best, alcanzándose con ellos típicamente resultados de second best. La razón es sencilla:
\begin{lnum}
    \item Si bien es cierto que el sistema permite que la información privada de cada Agente sea tácitamente revelada de manera veraz al Principal, esto solo se logra destinando recursos (rentas de información) para conseguirlo. Es decir, para lograr que cada tipo diga la verdad, el Principal debe dejarle una renta de información positiva a cada uno, lo que reduce su ganancia total. Entonces surge un dilema:
    \begin{la}
        \item El Principal puede conseguir la eficiencia asignando rentas informativas a cada tipo, lo que reduciría sus ganancias esperadas; o,
        \item Puede distorsionar los contratos, sacrificando algo de eficiencia en algunos contratos (en términos de bienestar general).
    \end{la}
    \item Como maximizador de beneficios, el Principal optará por la segunda opción, por lo que, en síntesis, el Principal modificará el nivel de esfuerzo, producción, calidad, o cualquier otra cosa, asignado a ciertos tipos para hacer menos atractivo su contrato a los otros tipos. Es decir, debe distorsionar las decisiones de ciertos tipos para, haciéndolas menos atractivas, “ahuyentar” a otros de hacerse pasar por ellos.\footnote{%
        \textit{No es por los varios miles de francos que tendrían que gastar para cubrir los vagones de 3ª clase o para tapizar los bancos que esta compañía tiene vagones descubiertos y bancos de madera; con gusto haría este sacrificio en aras de su popularidad. Su objetivo es evitar que el viajero que puede pagar el viaje en 2ª clase vaya a la 3ª clase. Perjudica a los pobres no porque quiera personalmente que sufran, sino para asustar a los ricos. […] En resumen, es también por este mismo motivo que las compañías, después de haberse mostrado casi crueles para los viajeros de 3ª clase, avaros para los de segunda, devienen pródigas para los de primera. Después de haber negado lo indispensable al pobre, se da lo superfluo al rico.} 
    
    (DUPUIT, 1849).
    } 
\end{lnum}

\textbf{NOTA.-} En general, se da en estos modelos una propiedad de “no distorsión por arriba” (no distortion at the top), por la cual para aquellos tipos por los que ningún otro tipo desea hacerse pasar (constituyen ese “top”, aunque según el modelo pueden ser los buenos o los malos desde la óptica del Principal) se alcanza la eficiencia, pero no para el resto de tipos, cuyas acciones exigidas por contrato serán ineficientes.

En conclusión, la asimetría informativa genera ineficiencias por la ventaja que da a algunos intervinientes.\footnote{%
    Por comparación con la situación de información simétrica, el Principal siempre pierde bienestar por dos vías: 
\begin{lnum}
    \item Debe garantizar una utilidad por encima de su nivel de reserva (U) a ciertos tipos: renta de información; y,
    \item Debe distorsionar las decisiones de algunos tipos para que no se hagan pasar por quienes no son.
\end{lnum}
Y respecto a los Agentes, la asimetría les reporta, a todos, al menos la misma utilidad que con información simétrica (la de reserva), haciendo que algunos tipos obtengan rentas de información por encima de dicha utilidad de reserva.
} 

\subsubsection{Competencia entre Principales: \textit{competitive screening}}
Hasta ahora solo se ha considerado un Principal único (monopolistic screening), pero pueden darse casos de múltiples Principales que compitan por unos mismos Agentes (competitive screening). Tres son las principales consecuencias:

\paragraph{Planteamiento del problema de optimización}
En competitive screening, el problema de asignación de contratos se formula buscando un menú de contratos que satisfaga (ROTSCHILD y STIGLITZ, 1976): (i) restricción de participación; (ii) restricción de compatibilidad de incentivos; y, (iii) restricción de beneficio nulo. 

Su desarrollo analítico ya no es un problema de optimización basado en Multiplicadores de Lagrange, sino un equilibrio de Nash mediante las técnicas de Teoría de Juegos, pues cada competidor deberá considerar la actuación desus competidores. En concreto, un contrato óptimo de este equilibrio será aquel que no solo sea aceptado por algún tipo de Agentes, sino que ningún otro Principal pueda ofrecer un contrato distinto que sea preferido por algún tipo de Agente y que le genere un beneficio esperado mayor.

Sin embargo, es común en la práctica solucionar el problema mediante la concepción de un planificador competitivo: esto es, resolver el problema de optimización introduciendo la restricción adicional de zero-profit para cada contrato ofrecido.\footnote{%
    Ya que hemos pasado a la situación en la que los Principales, a pesar de tener todo su poder de  negociación, no pueden obtener beneficios de ello.
} 

\paragraph{Tipos de Equilibrio: agrupador o separador}
Se pueden dar dos categorías de equilibrios:
\begin{la}
    \item Agrupador (pooling equilibria). Solo hay un contrato para todos los tipos de Agente; o,
    \item Separador (separating equilibria). Un contrato de equilibrio diferente para cada tipo de Agente. 
\end{la}

Obviamente, lo deseable para el Principal es poder separar o discriminar entre Agentes, como en el caso de Principal único, por lo que típicamente los equilibrios separadores le serán preferidos. Por desgracia, en estos problemas pueden darse casos en los que no existe equilibrio, ya sea de ninguna clase o solo de alguna, como en el modelo de ROTSCHILD y STIGLITZ (1976).\footnote{%
    En su análisis, ROTHSCHILD y STIGLITZ (1976) llegan a demostrar que en su modelo del mercado de seguros no existía equilibrio agrupador y la existencia de equilibrio separador no estaba asegurado (llamativa paradoja).
}

\paragraph{Análisis del Bienestar}
Respecto al bienestar, la distribución de ganancias entre Principales y Agentes pasa a estar sesgada hacia los Agentes. 
La presión competitiva entre Principales lleva a que los Agentes puedan exigir un pago que cubra su utilidad de reserva, su renta de información (en su caso) y ahora además una cuantía adicional ante la amenaza de elegir otro Principal que le reporte mayores pagos, por lo que los Agentes se reparten todo el resultado de la relación y los Principales tienen un beneficio nulo. 
Además, ahora no todos los Agentes ganan necesariamente con la información privada, pudiéndose dar el caso de que ciertos tipos reduzcan su bienestar (utilidad esperada) respecto a la situación de información simétrica. Esto es lo que abrirá la puerta a la siguiente categoría de modelos, la señalización.


\subsection{Ejemplo.- Modelo de ROTSCHILD-STIGLIZ}
Supongamos una Asegurador que se enfrenta a dos tipos de clientes, los de salud débil y los de salud fuerte. Cada cliente es consciente de qué tipo de salud tiene pero la Aseguradora, a priori, no. Por lo tanto, deberá ser capaz de diseñar:
\begin{la}
    \item Un contrato de forma que cada tipo de cliente se auto-selecciones por ser aquel tipo de contrato que más le beneficie; o,
    \item Un contrato que pondere las distribuciones poblacionales de manera que, independientemente del tipo de cliente que acceda a él, este sea aceptable para todas las partes (por ser también el único disponible).
\end{la}

\imagenfit{Modelo_ROSTCHILD_STIGLITZ_Seguros.png}{Plano de riqueza disponible en caso de contingencia ($w_a$) o no contingencia ($w_{na}$)}{Elaboración propia}

El Gráfico muestra la renta de todo cliente en caso de enfrentarse a una intervención que deba pagar. El Punto $E:(\bar w_{na}, \bar w_a)$ marca la intersección de la renta en caso de accidente con la renta en caso de no-accidente, siendo la bisectriz los puntos de equivalencia (accidente – no-accidente) y el Punto $A$ un posible punto óptimo para una compañía aseguradora, donde se cobre una prima sobre la renta disponible y se concede una compensación ($S=S-\varrho$) en caso de accidente. 

Ahora bien, la distribución de utilidad de cada Agente es diferente, pues los clientes sanos tendrán más utilidad al disponer de más renta sin considerar accidente que los Agentes no sanos, que tendrán más utilidad considerando la renta ex post (puesto que implícitamente valoran el riesgo de tener un accidente como alto). Es decir, los Agentes sanos ($H$) tienen preferencias cuasi-lineales respecto a la riqueza sin accidente ($w_{na}$) mientras que los agentes enfermos ($S$) tienen preferencias cuasi-lineales respecto a la riqueza en caso de accidente ($w_a$).

\imagenfit{Dist_Utilidad_TipoAgente.png}{Curvas de indiferencia en función del tipo de Agente: Preferencias cuasi-lineales en función del tipo de agente}{Elaboración propia}

Consideraremos tanto el Punto $A$ como el Punto $C$ como candidatos para un contrato \textit{pooling} (de distribución poblacional pura). 


\imagenfit{NO_POOLINGEq_ROTSCHILD_STIGLITZ.png}{Equilibrio separador en el modelo: Imposibilidad de realizar un equilibrio agrupador}{Elaboración propia}

El Gráfico de la izquierda muestra como, dada una distribución poblacional de tipo ($2/3$ de no-sanos frente a $1/3$ de clientes sanos). No podremos obtener un equilibrio agrupador, puesto que $\alpha$ es mejorable por cualquier contrato que se encuentre en la Región de $\beta$ (roja) ya que la utilidad que ofrece a los clientes sanos es mayor y estos optarán por éste. Como estos son precisamente los clientes que queremos atraer, otra Aseguradora captará a todos los clientes sanos ofreciendo un contrato ligeramente mejor, mientras que todos los no-sanos se quedaran en la del equilibrio agrupador. 

El único equilibrio estable en esta disposición es, por tanto, un resultado de segmentación de contratos como el que vemos en la imagen de la derecha ($\alpha_S,\alpha_H$). En este, el mejor resultado posible equivale a la Restricción de Compatibilidad de Incentivos, esto es, la Utilidad del contrato ofrecido a los clientes no-sanos es mayor o igual que la ofrecida a los clientes sanos, por lo que estos no tendrán incentivos para cambiar su contrato y manifestarse como sanos. Obtenemos así el Punto $\alpha_S$ y $\alpha_H$ y una prima y sistema de co-pago (o pago integral en caso de los clientes no-sanos) asociado a cada tipo de cliente. 

Sin embargo, en ROSTCHILD-STIGLITZ también indican que el equilibrio generado por un menú de contratos (auto-selección) es también inestable, puesto que si cambiaramos la distribución poblacional y la recta A-E quedara ahora por encima de la curva de indiferencia $H’$ (por ejemplo, 10\% insanos y 90\% sanos), entonces existiría un contrato óptimo (en términos de la Aseguradora y los Agentes) a modo de pooling tal que un solo contrato ofrecería mayor utilidad a todos los Agentes y mayores Beneficios a la Aseguradora, frente al contrato de segmentación que ofrece Beneficio nulo puesto que se encuentra sobre la recta de primas actuarialmente justas. 


\subsection{Modelos de selección adversa: signalling }
Los modelos de señalización son aquellos en los que el Agente es la parte activa en la mitigación de la asimetría informativa.\footnote{%
    El análisis pionero de esta familia de modelos se debe a SPENCE (1973), al demostrar que la educación podía servir como señal a las empresas (empleadores) de que una persona es simplemente capaz de aprender, independientemente de los contenidos que incluyan el título académico. Por lo tanto, esos contenidos pueden llegar a resultar irrelevantes, no aportando nada al desempeño laboral en la empresa: la educación puede ser una mera señal. Ahora bien, esto es un supuesto y no un resultado del análisis, por lo que este análisis no implica ni mucho menos defiende que la educación no aporte absolutamente nada, es solo un supuesto simplificador a fin de analizar de manera aislada el rol señalizador de la educación.
} En concreto, el Agente, después de conocer su tipo, pero antes de la firma del contrato, puede emitir una señal sobre su tipo que sea observada por el Principal. Los caracterizadores son los mismos que en selección adversa: 
\begin{lnum}
    \item Asimetría informativa ex ante, referida a una característica o información potencialmente heterogénea que solo conoce una de las partes (Agentes); y,
    \item No puede ser revelada verazmente sin coste.
\end{lnum}

¿Por qué querría un Agente mitigar la asimetría informativa? En la relación de agencia, como se ha visto, la información privada favorece en ocasiones a los Agentes, pero no siempre. En modelos de competitive screening, es posible que ciertos tipos de Agentes se vean perjudicados respecto a la situación de simetría informativa.\footnote{%
    Y, con más generalidad, también pueden darse muchos casos en los que los contratos de auto-selección no sean factibles o rentables para el Principal/es, por lo que una sencilla remuneración en función del tipo medio o esperado (como en el modelo de lemons) tenderá a favorecer a unos tipos y a perjudicar a otros.
} En consecuencia, estos tipos de Agentes perjudicados desearán erradicar la asimetría informativa y revelar su información privada a través de mecanismo que se denomina señal. 
Dicho de otro modo, solo tienen sentido estrategias de señalización en aquellos casos en los que ciertos tipos de Agentes salen perdiendo con la asimetría informativa.
Ahora bien, como en toda relación de agencia, la mera revelación de información no es creíble, pues ciertos tipos querrán falsearla, así que, para que sea efectiva la señalización debe ser a prueba de falseamiento. 
Una señal puede definirse como cualquier actividad o decisión que prueba que el Agente que la emite posee cierta habilidad, característica o información (es decir, su tipo). Para que sea informativa, solo los Agentes de ese tipo deberían estar interesados en emitirla, es decir, obtendrían un beneficio de revelar la información y ningún otro tipo querría imitarla, por lo que debe ser prohibitiva para otros tipos. 

\imagenfit{EsquemaTemporal_Señalizacion.png}{Esquema temporal de un problema de señalización}{Apuntes ICEX-CECO. Tema 3.A.13}
 
\subsubsection{Los elementos del problema}
Los elementos del problema son muy similares a los del problema de selección adversa, salvo algunas particularidades. Se parte de la misma relación entre Principal y Agente, aunque se supondrá un conjunto de Principales que compiten por contratar a los Agentes, para que lleve a cabo en su nombre alguna acción o tarea a cambio de una remuneración o contrapartida. En el caso de referencia que se venía usando, se suprime la variable esfuerzo, suponiendo que la asimetría surge simplemente alrededor de una característica innata del Agente: su productividad.

\begin{center}
\renewcommand{\arraystretch}{1.25}
\setlength{\tabcolsep}{6pt}
\begin{tabular}{|p{0.18\linewidth}|p{0.39\linewidth}|p{0.39\linewidth}|}
\hline
 & \textbf{Respecto del Agente} & \textbf{Respecto del Principal} \\
\hline
\textbf{¿Cómo son?} &
Heterogéneos y múltiples -- {\scriptsize Existe heterogeneidad y pluralidad de Agentes (cada uno con su tipo, denotado por $i$). Los Agentes difieren en su productividad (por simplicidad, se normaliza para que coincida con $i$).} &
Múltiples -- {\scriptsize Se considerará que hay varios Principales que compiten entre sí por contratar Agentes, por lo que la presión competitiva hará que el beneficio esperado de los Principales sea nulo.} \\
\hline
\textbf{¿Qué quieren?} &
Maximizadores de utilidad -- {\scriptsize La utilidad obtenida tras la firma del contrato viene definida por: (i) la remuneración según su productividad esperada, considerada como un bien; y, (ii) la desutilidad, un mal, de educarse dada su productividad $v_i(e)$. Para el modelo base, se asume separabilidad: $U_i(w_i,e)=u(w_i)-v_i(e)$.}\footnote{%
    La utilidad de reserva, $\bar U$, ya no juega un papel relevante dada la competencia de Principales.
} &
Maximizadores de beneficio -- {\scriptsize Aunque este será nulo debido a la competencia entre Principales.} \\
\hline
\textbf{¿Qué información poseen?} &
Información privada -- {\scriptsize En el ejemplo considerado, esta información corresponde a su productividad (asimilada al tipo $i$) y, en consecuencia, la desutilidad o coste que les genera educarse.} &
Distribución de probabilidad de los tipos -- {\scriptsize El Principal conoce la distribución de tipos/productividad $F(i)$, que describe cómo se reparte la probabilidad de cada tipo (y su productividad asociada) en términos poblacionales.} \\
\hline
\textbf{¿Qué eligen?} &
Revelar información -- {\scriptsize Antes de aceptar el contrato, los Agentes pueden emitir una señal para diferenciarse ante el Principal. Se consideran señales puras: $e$ denota el nivel de educación elegido como señal, y cada tipo presenta un coste distinto $v_i(e)$.}\footnote{%
    Esto es, actividades que no proporcionan ninguna utilidad o beneficio al Agente más allá de señalizar. Aunque este supuesto puede resultar muy restrictivo, se hace para aislar la actividad señalizadora de otras motivaciones. Además, permitir en ciertos casos que la señal tenga una utilidad o beneficio para el Agente no modifica cualitativamente los resultados del análisis.
} & 
Remuneración vinculada a cada señal -- {\scriptsize Tras observar la señal, el Principal actualiza racionalmente sus creencias y obtiene la distribución condicional $F[i\mid e]$. En función de ella, diseña una remuneración asociada a cada mensaje y ofrece el contrato correspondiente al Agente.} \\
\hline
\end{tabular}
\end{center}

\subsubsection{Planteamiento del problema: Equilibrio dinámico}
El planteamiento riguroso de un problema de señalización es complejo, pues es un problema dinámico (en etapas) y su solución es un equilibrio que requiere consistencia tanto de las decisiones de Agentes (señales) como de Principales (creencias y contratos). Por lo tanto, no se puede resolver mediante un programa de optimización con Restricción de Participación y Restricción Presupuestaria (tampoco si añadiésemos la Restricción de Beneficios), sino que requiere herramientas avanzadas de teoría de juegos. 

Estos equilibrios (Equilibrios Bayes-Nash) se definen mediante dos componentes, un perfil de estrategias de los intervinientes y el sistema de creencias o conjeturas de la parte no informada, siendo estas creencias el componente que plantea ciertas dificultades técnicas.

\subsubsection{Equilibrio: Separadores y Agrupadores}
Como en los modelos de selección adversa, sin pérdida de generalidad podemos decir que se pueden dar dos clases de equilibrios, pero en este caso será función de la señal y no del tipo de agente (que se desconoce al igual que el problema anterior).

\paragraph{Equilibrios agrupadores (pooling)}
Si todos los tipos de Agentes emitiesen la misma señal, lo más razonable es que ningún Agente incurra en costes de señalización: $e(i)=0\;\;\forall i$. Ante esta señal, un Principal solo puede ofrecer a todos los Agentes el mismo contrato o remuneración determinada por la esperanza ex ante de tipos (parte izq.) que haga sus beneficios esperados nulos (parte dcha.):

$$w_A=w(\mathbb E[i])\Longrightarrow\mathbb E[\Pi(i)]=w(\mathbb E[i])$$

En ciertos casos la existencia de esta clase de equilibrios no está garantizada, y cuando sí lo está puede que no resulten razonables, ya que ciertos Agentes preferirán señalizar.\footnote{%
    En efecto, es posible que algunos agentes mejorasen respecto al equilibrio agrupador si señalizan su tipo, por lo que tendrían incentivo a desviarse y por lógica no podría tratarse de un equilibrio. Este carácter poco razonable de ciertos equilibrios se debe a unas creencias de los Principales que pueden resultar miopes, de ahí que se suela aplicar el denominado “criterio intuitivo” de CHO y KREPS (1987) para descartar tales situaciones.
}

\paragraph{Equilibrios separadores (separating)}
Si cada tipo de Agente emite una señal distinta y el Principal puede así identificar a cada tipo de Agente, ofreciéndole una remuneración acorde a su tipo revelado verazmente, se consigue minimizar la asimetría informativa. En el ejemplo seguido, el nivel de educación elegido por cada Agente será $e(i)=(v_i(e))^*,\;\;\forall i$. El Principal actualizará sus creencias a priori y ofrecerá una remuneración de acuerdo con el tipo,  $w^s_i=\Pi(i)$.

Pero ¿cómo debe ser esa señal para lograr esto? Lógicamente, para un Agente de un cierto tipo $i$ la señal debe ser lo suficientemente costosa para evitar que sea imitada por otros tipos, pero lo suficientemente barata para que al propio Agente le compense emitirla, por comparación no solo con lo que obtendría haciéndose pasar por otro tipo, sino también comparando con un equilibrio agrupador (de lo contrario, no sería racional emitir señal y no existiría equilibrio separador). Esto es, el Agente que emite señal debe cumplir:

$$U^s_i(w_i^s,e^*)\ge U_i^A(w^A,0)$$
  
Que esta desigualdad se cumpla puede depender de diversos factores clave, como la distribución de tipos (su media, su dispersión…) y los costes relativos o desutilidad relativa de emitir señal para cada tipo.

\subsubsection{Análisis de Bienestar}
Los Principales tendrán un beneficio nulo en todos los equilibrios, así que quedarán indiferentes entre un equilibrio y otro. La situación de los Agentes puede ser más compleja y dependerá del tipo de equilibrio.

\paragraph{Bienestar de los Agentes en equilibrio respecto a información simétrica}
Respecto a la situación de simetría informativa:
\begin{lnum}
    \item En un equilibrio agrupador con señal nula en que se remunera uniformemente al nivel del tipo medio o esperado:
    \begin{la}
        \item Beneficiados: Aquéllos cuyo tipo, de ser conocido, les proporcionaría una remuneración inferior a esta remuneración media. 
        \item Perjudicados: Los que poseen tipos a los que les correspondería una remuneración superior a la media. 
    \end{la}
    \item En el caso de equilibrio separador:
    \begin{la}
        \item Beneficiados: Aquellos que emitan una señal nula (y por tanto sin incurrir en este coste) y que cobren conforme a su tipo, se encontrarán igual que con información simétrica;
        \item Perjudicados: Los Agentes que deban emitir una señal costosa para poder ser remunerados conforme a su tipo real se verán perjudicados exactamente en la cuantía del coste de emitir la señal.
    \end{la}
\end{lnum}

\paragraph{Bienestar de los Agentes entre equilibrios}
Como se esperaría, los tipos de Agentes que en el equilibrio separador reciban un pago neto de costes de señal superior al pago del equilibrio agrupador, estarán mejor en dicho equilibrio separador. Sin embargo, también habrá Agentes que estén en el caso opuesto y prefieran en equilibrio agrupador. Lo crucial aquí es que los Agentes que deban emitir señal para diferenciarse ganen bienestar emitiéndola, pues de lo contrario, como se indicaba más arriba, no sería razonable un equilibrio separador.

En resumen, se alcanzan conclusiones muy similares al caso de selección adversa o screening. La solución de señalización, aunque permita revelar tácitamente la información privada de los Agentes, no supone un first best, sino que se trataría de un second best puesto que, en el caso de equilibrio separador, igual que en selección adversa, los Agentes deben dedicar recursos a mitigar la asimetría informativa (en este caso, a través de señales costosas), constituyendo la señal un medio costoso para redistribuir renta de unos Agentes a otros. Y en caso de equilibrio agrupador (sin coste de señal), el análisis de bienestar puede resultar ambiguo, pudiendo darse situaciones en las que las decisiones de producción sean eficientes, aunque la distribución de ganancias no sea la obtenida con información simétrica.

\subsubsection{Ejemplo.- Educación (SPENCE, 1973)}
La señalización (signalling) propuesta formalmente por primera vez por SPENCE en su artículo \textit{Job Market Signalling} (1973) consiste en la transmisión creíble de información privada desde las partes informadas a las partes desinformadas haciendo uso de una tecnología relativamente más costosa para los agentes con un bien de calidad baja.

Es precisamente ese mayor coste para los “malos” agentes lo que permite que la señal sea creíble para la parte no informada ya que, si enviar una señal a la parte informada no tuviese coste alguno, todos los agentes informados tratarían de informar a los desinformados de que su calidad es alta y por supuesto, estos últimos no otorgarían racionalmente ningún valor a esa señal. Es decir, el primer supuesto del modelo es que existe una correlación negativa entre el coste de señalizar y aquella característica que es apreciada como positiva para la empresa: the costs of signaling are negatively correlated with productive capabilities. Existen dos tipos de Agentes atendiendo a su productividad:
\begin{la}
    \item Productividad alta: A; y,
    \item Productividad baja: B
\end{la}

Pero la productividad del trabajador no es observable por la empresa hasta después de haberlo contratado. Es decir, tenemos un marco en el que obtienen un ingreso igual a la productividad individual y no pueden verificar la productividad ex-ante. Por lo que sucedería lo mismo que en el ejemplo de AKERLOF, los trabajadores más productivos son expulsados del mercado, ya que el salario que ofrecen las empresas cae ante la posibilidad de que el trabajador sea de productividad baja (ponderación del peso o proporción poblaciónal de cada agente por PMg de cada agente, que en equilibrio competitivo sería igual al salario): este equilibrio daría una posición ventajoso para los “malos” y una desventajosa para los “buenos”.
Para distinguirse, los trabajadores usarán la educación a modo de señal para la empresa ya que ésta depende exclusivamente de su decisión individual y se ve condicionada por: (i) el grado de producividad implícito del Agente; y, (ii) el grado de educación obtenido / bien visto por el Principal. El coste de alcanzar ese nivel de educación es menor para los individuos de productividad alta (A) que para los individuos menos productivos (B), de forma que:

$$v_A(e^*)<v_B(e^*)$$
 
De este modo, los individuos estarán interesados en incurrir aquel nivel educativo que pueda servir como señal creíble, sabiendo que:
\begin{la}
    \item Si $E\ge \widehat E$, siendo ésta última el grado de educación óptimo, el empresario pagará el salario correspondiente a la Alta productividad ($PMg_A=w_A$); y,
    \item Si es inferior, el salario asignado será el de Baja productividad ($PMg_B=w_B$).
\end{la}

Para plantear gráficamente este problema solo debemos identificar el siguiente conjunto de restricciones, a partir del cual obtendremos la cantidad óptima de señal (eduación) que el Agente y el Principal estarán dispuestos a ofrecer o solicitar, respectivamente:
\begin{lnum}
    \item $U(w_A,\widehat E)=w_A-v_A(\widehat E)>\bar U$. Restricción de Participación; y,
    \item $w_B<w_A-v_A(\widehat E)$ y $w_A-v_B(\widehat E)<w_B$. Restricción de Compatibilidad de Incentivos;
\end{lnum}

La decisión de los trabajadores concierne el grado de educación que maximizará su beneficio, teniendo en cuenta que la educación es costosa pero puede permitirles obtener un salario mayor: ¿cuántos años de educación adquirir? Atendiendo al mercado en general, la pregunta relevante es ¿todos los trabajadores elegirán la misma educación, o niveles diferentes? ¿el nivel elegido depende de la productividad? La respuesta a estas preguntas dependerá del CMgE relativo de la educación para trabajadores con productividad alta, y del salario ofrecido por las empresas. Existe así, para los trabajadores, un evidente trade-off entre salario recibido y educación obtenida. La decisión de las empresas concierne las combinaciones de salario y educación que ofrecen a los trabajadores. ¿Cuánto salario deben ofrecer respecto a la educación observada? ¿Deben ofrecer el mismo salario para todos? ¿o diferentes salarios para cada nivel de educación? ¿a cuánto debe ascender el salario? Responder a estas preguntas y las relativas a la decisión del trabajador consiste en caracterizar los posibles equilibrios que puedan aparecer.

En este contexto, un equilibrio será un conjunto de combinaciones salario–educación para los que las empresas obtienen beneficios nulos, ninguna empresa puede mejorar su beneficio y los trabajadores maximizan su ingreso individual. Existen generalmente dos tipos de equilibrios en este contexto:
\begin{la}
    \item Equilibrios separadores (separating equilibria) en los que los trabajadores eligen diferentes niveles de educación en función de su productividad; y,
    \item Equilibrios agrupadores (pooling equilibria) en los que los todos los trabajadores eligen el mismo nivel de educación independientemente de su productividad.\footnote{%
        En el modelo de SPENCE, si la proporción de trabajadores de productividad alta es muy elevada, un equilibrio aunador (otorgar un mismo salario a todos los trabajadores, independientemente de su productividad) podría llevar a todos los trabajadores a un mayor nivel de utilidad que el equilibrio separador (otorgar un salario diferente en función de la productividad del trabajador).
    }
\end{la}

Precisamente como hay que incurrir en un coste adicional (educación) se trata de una solución de second best y no se cumple el Primer Teorema Fundamental de la Economía del Bienestar.\footnote{%
    El propio SPENCE muestra empíricamente que los trabajadores tienen una tendencia a sobreeducarse con respecto a lo que sería eficiente. Algunos autores han utilizado el modelo de señalización de SPENCE para explicar las consecuencias de la implantación de los planes Bolonia en España. En la medida en que, dicen, los actuales planes de Bolonia son más fáciles que las antiguas licenciaturas (por el mero hecho, siquiera, de que se da menos contenido al ser las carreras más cortas), los costes del esfuerzo son menores que antes. De esta manera, al nivel de educación exigido anteriormente por las empresas (digamos, una carrera, $\widehat E$) se produce un equilibrio agrupador, pues al ser las carreras más fáciles (i.e. menos costosas), $\widehat E$ se encuentra en un punto más a la izquierda y los agentes tienen más incentivos a obtener ese nivel de educación. De ahí el fenómeno que hemos observado recientemente de que las empresas están exigiendo, además de una carrera un máster (pues es necesario exigir un mayor nivel de educación para que se genere un equilibrio separador).
}

\subsection{Selección adversa y riesgo moral: complementarierdad}















 












\clearpage
\section{Conclusión} 

\subsection{Extensiones y relación con otras partes del Temario}


\subsection{Opinión}



\subsubsection{Idea final (Salida o cierra)}


\clearpage
\pagenumbering{gobble}
\MyBibliography
\MyBibItem{Autor}{Título}{Año}{DOI -- LINK}




%ANEXOS
\clearpage
\anexo{\textbf{Preguntas Test}}
%\input{\TemaCode/questions.tex} 




\end{document}